\documentclass[10pt,a4paper,twoside,bg=print,twocolumn,openany,nodeprecatedcode]{dndarticle}
\usepackage[utf8]{inputenc}
\usepackage{tabulary}
\usepackage[normalem]{ulem}
\usepackage{color,soul}
\usepackage{float}
\usepackage{caption}
\usepackage{wrapfig}
\usepackage{tocloft}
\usepackage[toc]{multitoc}
\usepackage[hidelinks]{hyperref}
\raggedbottom
\extrafloats{2000}
\cftsetindents{section}{0em}{0em}
\cftsetindents{subsection}{0em}{0em}
\cftsetindents{subsubsection}{0.2em}{0em}
\setcounter{tocdepth}{1}
\renewcommand*{\multicolumntoc}{3}
\setlist[description]{nolistsep,listparindent=3pt,topsep=6pt}
\date{} 
\begin{document}
\part{Affair on the Concordant Express}
\textit{An Adventure for 9th-Level Characters}

Wanted for crimes across the multiverse, an outlaw known as \textbf{the Stranger} is currently a prisoner aboard the Concordant Express, an interplanar train destined for Mechanus, where the outlaw is set to stand trial. In this heist, the characters must obtain a list of names \textbf{the Stranger} has committed to memory. The names in question are the true names of several powerful demons, devils, and yugoloths.

\section{Adventure Background}
The interplanar outlaw known as \textbf{the Stranger} had humble origins. \textbf{The Stranger} was born into a slum and orphaned as an infant. Early on, they used their knack for magic to masquerade as a lost noble scion and earn patronage as a spellcasting student. Inevitably, the school \textbf{the Stranger} was attending discovered the ruse, but \textbf{the Stranger} always slipped away and reemerged at an even more prominent academy in another city with a new false identity.

\textbf{The Stranger} perpetrated many crimes, none greater than the attempted theft of a book that held the true names of many demons, devils, and yugoloths (see "True Names" below). When \textbf{the Stranger} tried to flee with the book, a fiendish ward caused the book to disappear. By then, \textbf{the Stranger} had already committed many of the true names to memory. \textbf{The Stranger} then set off across the multiverse, making deals with various Fiends in exchange for keeping their true names secret. \textbf{The Stranger} survived all these dealings unscathed for years, but their luck eventually ran out.

Fiendish foes set their sights on \textbf{the Stranger}, forcing \textbf{the Stranger} to flee to the Material Plane and take refuge in the city-state of Akharin Sangar, which is ruled by a powerful angelic being. (Akharin Sangar is detailed further in Journeys through the Radiant Citadel). But \textbf{the Stranger} underestimated the vigilance of Akharin Sangar's celestial host and was caught impersonating an angel. When \textbf{the Stranger} refused to divulge any information about themself, a deva enforcer named Omid was tasked with delivering \textbf{the Stranger} to Mechanus to stand trial. The prisoner was clamped in \textit{dimensional shackles} and taken aboard the Concordant Express for safe transport.

\subsection{True Names}
Every Fiend has a true name that it tries to keep secret. One who knows a Fiend's true name can use this name to bind that Fiend to service. Other powerful creatures such as Celestials and Fey might have true names that function similarly.

True names aren't restricted to spoken language or script. A creature's true name could be a specific gesture, sound, or other type of expression, such as the tolling of an iron bell, a mathematical equation, or a sequence of notes played on a particular kind of instrument.

\section{Adventure Hooks}
If you're using the Golden Vault as the hook for the adventure, skip ahead to the "Using the Golden Vault" section. Otherwise, choose one of the following adventure hooks.

\subsection{Blood War Balance}
A planar faction known as the Fixers aims to restore balance to the Blood War (the eternal feud between demons and devils) by obtaining the true names of a pit fiend commonly known as Karnyros, a balor called Errtok, and a marilith known as Hexalanthe. A representative of the Fixers hires the characters to obtain these Fiends' true names from \textbf{the Stranger}.

\subsection{Fiendish Extortion}
A high-ranking member of the Fated, a faction in Sigil, has been captured and is being ransomed by an ultroloth known as Zeevok. Faction leaders hire the characters to obtain Zeevok's true name so the Fated can force Zeevok to release his captive without collecting a ransom.

\subsection{Personal Reasons}
The characters seek the true name of a Fiend that has bedeviled them in the past.

\section{Using the Golden Vault}
If you're using the Golden Vault as a patron, a golden key is delivered to the characters in whatever manner you deem fit. When the characters use this key to open their music box, the lid pops open and a soothing voice says the following:

\begin{DndReadAloud}{}
"Greetings, operatives. The Golden Vault needs you to obtain the true names of three powerful denizens of the Lower Planes: the pit fiend Karnyros, the balor Errtok, and the marilith Hexalanthe. Their true names are known to an outlaw called \textbf{the Stranger}, who is being transported to Mechanus on an interplanar train called the Concordant Express. This quest, should you choose to undertake it, requires you board the train, locate \textbf{the Stranger}, and acquire the true names from them. Expect to be contacted shortly by a quadrone named Glitch, a former train operative who now works for us. Good luck, operatives."

\end{DndReadAloud}

Closing the music box causes the golden key to vanish.

\section{Glitch the Quadrone}
Not long after agreeing to undertake the adventure, the characters encounter a friendly \textbf{quadrone} named Glitch. (When and where this happens is up to you.) The quadrone used to oversee hospitality on the Concordant Express. For this reason, it was imbued with the ability to speak Celestial, Common, and Infernal in addition to Modron.

Conflicting directives from two pentadrone superiors caused a glitch that prompted the confused quadrone to quit its job. Glitch has become a free-thinking being with no desire to reintegrate into the modron hierarchy.

Eager to help the characters, Glitch speaks to them in Common:

\begin{DndReadAloud}{}
"Ah, hello! Yes, good, finally—friends looking to catch \textbf{the Stranger}—\textbf{the Stranger} who has already been caught. Here, take this map. It shows the train cars. Here are your tickets and a pen to write down names. Not ordinary names—secret names! I am Glitch, by the way."

\end{DndReadAloud}

Glitch hands the characters an old blueprint of the Concordant Express (see map 11.1), one train ticket per character, and an ink pen. The tickets are made from paper-thin sheets of brass and grant passage to Mechanus aboard the Concordant Express. Each ticket is stamped with a boarding time, and the characters have 24 hours until the train departs.

Glitch takes the characters to the train when they're ready (see "Catching the Train" below) and stays with them in the interim. Glitch also shares the following information:


\begin{description}
\item[Configuration.]
Like most trains, the Concordant Express includes an Engine Car at the front and a Caboose at the back. The number of other train cars and their configuration can change, so Glitch can't say much more about the size of the train or the order in which the cars will be arranged. The configuration of the cars might be different than what's shown on the blueprint. (In fact, it is. See the "Train Configuration" section for details.)
\item[Fuel.]
The train is fueled by treasure, which is kept in the Engine Car.
\item[Location of the Stranger.]
\textbf{The Stranger} is likely being held in the Jail Car, which has a skylight.
\item[Patrols.]
Modrons patrol the exterior of the train.
\item[Teleportation Ward.]
Creatures can't teleport into or out of the train.
\end{description}

\section{The Concordant Express}
The Concordant Express is a clockwork locomotive that traverses the multiverse. Modrons constructed the train with mathematical precision, and they operate it in the same way. The train's conductor is its sapient Engine Car, which delivers passengers and cargo to their intended destinations on schedule.

No terrain is too rugged for the train. Like the landmasses that make up Mechanus, the wheels of the Concordant Express are a network of interlacing cogs. A series of mechanical arms rapidly places levitating tracks before the train, while a similar set of arms beneath the Caboose collects the trailing tracks and delivers them to the Engine Car. This system allows the train to traverse inhospitable environments and to adjust course as needed.

\subsection{Train Configuration}
The train has three mandatory cars and three optional cars, as described below.

\subsubsection{Mandatory Cars}
In this adventure, the Concordant Express has one of each of the following cars:


\begin{itemize}
\item The Caboose (area E1) is the rearmost car.
\item The Engine Car (area E9) is the frontmost car.
\item The train has one Jail Car (area E8), which is adjacent to the Engine Car (not the Caboose, as shown on map 11.1).
\end{itemize}

\subsubsection{Optional Cars}
Map 11.2 includes maps for six additional cars: Abacus Car (area E2), Aquarium Car (area E3), Cargo Car (area E4), Planartarium Car (area E5), Passenger Car (area E6), and Temple Car (area E7).

Pick any three of these cars and add them to the train in whatever configuration suits you, provided the mandatory cars stay where they are. (You can choose a different set of optional cars each time you run the adventure.)

You can add more than three optional cars to the train, at the risk of increasing the heist's length and difficulty. Conversely, you may opt for fewer cars for a shorter, easier heist.

\subsection{Catching the Train}
Embedded in Glitch's body is an \textit{amulet of the planes} that turns to dust (along with the rest of the quadrone) if Glitch is reduced to 0 hit points. Glitch uses this built-in amulet to cast \textit{plane shift}, delivering itself and the characters to a field in \textit{the Outlands} (see the \textit{Dungeon Master's Guide} for more information about this plane of existence). There they must wait for the train to arrive.

As the train approaches the characters' location, read the following text:

\begin{DndReadAloud}{}
You pass the time on a grassy plain, waiting for the train to arrive. A few miles away, a colossal spire of rock rises into the sky, its peak lost in the clouds.

Glitch grows fidgety and impatient. Suddenly, the distant blare of a horn draws Glitch's attention—and yours—to a column of smoke on the horizon. "Good," says Glitch. "Here it comes, right on time!"

A locomotive hurtles across the sky toward you, laying down magical tracks before it. As it draws near, you can see clockwork limbs disassembling the tracks behind the train and passing them forward so the locomotive can lay down new tracks in front of itself.

As the train descends to ground level, its horn blares once more. It swerves away from you, coming to a gradual stop until its caboose is a short distance away, pointing in your direction.

"Got your tickets?" asks Glitch. "Excellent. Now go. Show your tickets to the modron in the caboose. This is where I leave you. Farewell, friends!"

\end{DndReadAloud}

Glitch exits the adventure at this point, using its \textit{amulet of the planes} to return whence it came.

\subsubsection{Boarding the Train}
The Caboose (area E1) doubles as a train station. Inside it is a booth containing a \textbf{duodrone} attendant with a built-in ticket puncher. All passengers must get their tickets punched by this duodrone before they're permitted to explore the rest of the train.

If one or more characters refuse to get their tickets punched, the duodrone whistles (no action required), causing four tridrones to emerge from hidden compartments in the walls and act as described in area E1. If the characters then comply to avoid an altercation, the tridrones return to their compartments.

\subsubsection{Choo, Choo!}
Once the characters' tickets are punched and the train is ready to go, read the following text:

\begin{DndReadAloud}{}
Machinery whirs and clicks as the train shudders and creeps forward, gradually picking up speed. Another horn blast signals that the Concordant Express is once again on the move.

\end{DndReadAloud}

\subsection{General Features}
The Concordant Express has the following features:


\begin{description}
\item[Ceilings.]
Car ceilings are 15 feet high unless the text states otherwise.
\item[Doors.]
At each end of every car is a sliding metal door. The doors to the Jail Car are locked (see area E8 for details); all others are unlocked.
\item[Flight.]
The train has a flying speed of 60 feet and can't hover.
\item[Lighting.]
All cars are brightly lit by magical orbs embedded in the ceilings.
\item[Modrons.]
All modrons serving aboard the Concordant Express speak no languages other than Modron unless the text states otherwise.
\item[Teleportation Ward.]
Creatures can't teleport into or out of the train, or to any space within 120 feet of it. Any attempt to do so is wasted.
\end{description}

\subsection{Moving on the Train}
To reach \textbf{the Stranger}, the characters must make their way from the Caboose (area E1) to the Jail Car (area E8).

\subsubsection{Crossing between Cars}
Creatures can safely move from one car to another by stepping across couplers situated between the cars' exterior platforms.

\subsubsection{Falling off the Train}
If an effect would result in a creature falling off the train, the creature must succeed on a DC 14 Dexterity saving throw. On a success, the creature instead falls prone in the nearest unoccupied space on the train. On a failure, the creature is sucked into the train's undercarriage, taking 33 (6d10) bludgeoning damage and landing prone in an unoccupied space at either end of the nearest train car.

If a character falls off the Caboose, the train's track-retrieval arms snatch up the character and deposit them into an empty cell in the Jail Car (area E8).

\subsubsection{Guard Patrols}
Hostile modrons patrol the rooftops of the cars, attacking anyone they see who isn't supposed to be there. Whenever one or more characters reach the midpoint of a rooftop on any car other than the Jail Car, roll a d6 and use the Guard Patrols table to determine if an encounter occurs. Assume the patrol begins the encounter on the rooftop of an adjoining car. All the modrons have truesight, and characters on the rooftop have nowhere to hide.

No random encounters occur on the roof of the Jail Car until after the squad of pentadrones stationed there is defeated (see area E8 for details).

\begin{DndTable}[header=Guard Patrols]{cX}

d6 & Encounter\\
1–3 & No encounter\\
4–5 & Three tridrones and one \textbf{quadrone} leader\\
6 & Four quadrones and one \textbf{pentadrone} leader\\
\end{DndTable}

\subsubsection{Strong Winds}
Creatures climbing on the outside the train or flying within 15 feet of it while it's in motion are subject to strong winds, the effects of which are as follows:


\begin{description}
\item[Difficult Terrain.]
The wind-whipped exterior of the train is \textit{difficult terrain}.
\item[Ranged Attacks.]
Ranged weapon attacks are made with disadvantage.
\end{description}

\subsection{Planar Effects}
The Concordant Express moves across the known planes of existence by passing through invisible portals only the train's sapient Engine Car can find. The train's next stop is Mechanus, but it will pass through several other Outer Planes on its way there. While the train shields passengers from the harmful planes through which it travels, creatures aboard the train aren't entirely immune to the planes' effects.

The Planar Effects table contains effects for various planes of existence on the train's current route. Unless otherwise noted, each effect lasts until the train leaves the plane. Some effects apply only to creatures outside the train.

The train enters a new plane whenever you like. Roll or choose from the options on the table to determine which plane the train is on and its accompanying effect. Once a train leaves a plane, it doesn't return to that plane for the rest of the adventure.

\begin{DndTable}[header=Planar Effects]{cX}

d6 & Plane and Effect\\
1 & Acheron. The train flies over a vast battlefield where legions of devils clash with hordes of demons. Whenever one creature aboard the train deals damage to another creature aboard the train, both take 4 (1d8) psychic damage.\\
2 & Elysium. All creatures aboard the train feel a sense of inner peace. No creature aboard the train can make an attack or cast a spell that deals damage.\\
3 & Mount Celestia. Good-aligned creatures aboard the train gain a benefit of a \textit{bless} spell that lasts until the train leaves Mount Celestia.\\
4 & The Nine Hells. On one of the blisteringly cold layers of the Nine Hells (either Cania or Stygia), the train passes through a frigid gorge filled with whirlwinds of ice needles. Any creature outside the train takes 5 (2d4) cold damage plus 5 (2d4) piercing damage at the start of each of its turns. Modrons that normally patrol the outside of the train lock themselves in compartments for safety until the train leaves the plane.\\
5 & Pandemonium. Howling darkness surrounds the train. Creatures outside the train are deafened. Outside the train, bright light becomes dim light, and dim light becomes darkness.\\
6 & Ysgard. Silent, spectral warriors from Valhalla fly around the train on ghostly chariots pulled by ghostly steeds. Any creature aboard the train that would be reduced to 0 hit points drops to 1 hit point instead.\\
\end{DndTable}

\section{Train Cars}
The various train cars are shown on map 11.2. The earlier "Train Configuration" section explains how to choose and arrange the cars.

\subsection{E1: Caboose}
The characters are assumed to be in this car when the train gets underway. As the train begins to pick up speed and lifts into the air, a disembodied voice says, "Next stop: Mechanus!" in Common, Modron, Celestial, Infernal, Abyssal, and Primordial.

\subsubsection{Modrons in the Caboose}
The \textbf{duodrone} attendant who punched the characters' tickets has another function aboard the train: making minor repairs while the train is in motion. As the train gets underway, the duodrone dons gear from a nearby locker (see "Locker" below). It then departs to perform its duties, whistling as it marches to the next car.

Hidden in the walls of the Caboose are four tridrones that emerge from their compartments to attack passengers who refuse to get their tickets punched or who perform an unlawful act, such as threatening the duodrone or plundering the locker.

\subsubsection{Locker}
The locker contains a blue engineer's hat and a crew uniform consisting of a tool belt and a pair of grimy overalls labeled with a number.

\subsubsection{Rear Platform}
A railing on the car's rearmost platform prevents creatures from stumbling over the edge and into the dizzying blur of mechanical arms that retrieve the trailing railway tracks. A bronze ladder welded to the car provides access to its roof.

\subsection{E2: Abacus Car}
\begin{DndReadAloud}{}
Ten spherical modrons are lined up against the wall opposite a podium, behind which is a chalkboard covered with mathematical equations. Next to the podium stands a boxy modron wearing a wizard's cap. It taps the podium with a brass wand to get your attention. "Well met!" it says in Common. "I am the Math Wizard, Captain of Calculus, Admiral of Abaci, Unifier of Units! Welcome to my Sanctum Numerum, where any problem can be solved. I beseech thee, travelers, what numerical enigmas plague thee?"

\end{DndReadAloud}

The eleven modrons stationed here act as a mathematical calculator for passengers.

The self-proclaimed Math Wizard is a friendly \textbf{duodrone}. The Math Wizard speaks both Common and Modron.

The ten monodrones speak Modron only and are indifferent toward passengers. Each one holds a stack of ten placards numbered 0 through 9. When a creature inside the Abacus Car asks a mathematical question to which the answer is a number, the Math Wizard taps the podium with its wand and says, "Calculate!" In response, the monodrones select from their placards, shuffle around, and hold up their numerals to display the answer. The Math Wizard cheers "Huzzah!" after every successful calculation.

The monodrones do their best to calculate accurately, but sometimes they must estimate. Furthermore, if the ten monodrones are asked a question to which the answer exceeds 9,999,999,999 or the solution is uncountable, they scramble in circles while screaming in existential terror for a few moments before calmly returning to their positions.

\subsubsection{Chalkboard}
The chalkboard is covered with equations explaining, rather incomprehensibly, why the Great Modron March happens only once every 289 years. To anyone other than a modron, most of the equations are utter nonsense.

A character who studies the equations for at least 10 minutes can make a DC 25 Intelligence (Arcana) check at the end of that time. If the check fails, the character takes 3 (1d6) psychic damage and suffers a mild headache for the next hour. If the check succeeds, the character discovers a profound equation nestled among the others. Like a key, this equation unlocks something in the character's brain, with the effect determined by rolling on the Mysterious Equation table. A character who experiences this revelation gains no further benefit from studying the chalkboard equations. If any of the equations are erased, the modrons in the car turn hostile and attack, and no one can benefit from studying the equations.

\begin{DndTable}[header=Mysterious Equation]{cX}

d8 & Effect\\
1–3 & You gain \textit{inspiration}.\\
4–5 & You gain proficiency in one tool of your choice.\\
6 & You gain proficiency in one of the following skills of your choice: Deception, History, Insight, Investigation, Medicine, or Performance.\\
7 & You learn a new language of your choice.\\
8 & Your Intelligence score permanently increases by 2 (to a maximum of 20).\\
\end{DndTable}

\subsection{E3: Aquarium Car}
\begin{DndReadAloud}{}
This car is filled with clear water that isn't displaced when the door opens. The car's interior space is a large aquarium, complete with a coral reef and algae-coated ruins. Tiny fish dart through the ruins and circle an alabaster statue carved in the likeness of an angel clutching a decanter. Nearby fish scatter at the sight of you.

\end{DndReadAloud}

This car is completely flooded, and magic keeps the water from being displaced when either door is opened. Creatures that can't breathe underwater must hold their breath in the aquarium. Magic instantly dries creatures as they exit the water.

The aquarium can be slowly drained by removing a fist-sized plug near the base of the car's portside wall. As an action, a character within reach of the plug can try to pull it out, doing so with a successful DC 17 Strength check. Once the plug is removed, the aquarium's water level decreases at a rate of 1 foot per minute. It takes 15 minutes for all the water to drain from the car. The car's fish can't survive more than a few minutes without water.

\subsubsection{Angelic Statue}
The alabaster statue is a petrified \textbf{deva} with a \textit{decanter of endless water}. If any harm befalls the fish in the aquarium, the deva immediately reverts to flesh and chastises whoever is responsible. If the aquarium's plug was removed, the deva replaces it. If any water was drained from the car, the deva uses the decanter to replenish it. If any fish died, the deva uses \textit{raise dead} spells to revive one fish per day for ten days. Once all is set right, the deva flies from the train and casts \textit{plane shift} to return to Mount Celestia, taking the decanter with it.

\subsection{E4: Cargo Car}
\begin{DndReadAloud}{}
A battered war machine rests at one end of this spacious car. Acid drips from a hose-like weapon mounted on the front of the vehicle, leaving a sizzling hole in the train car's floor.

\end{DndReadAloud}

The war machine comes from the wastelands of Avernus, the first layer of the Nine Hells. The vehicle isn't operational, but its hose-like sprayer is. A creature standing within 5 feet of the nozzle can take an action to cause the weapon to spray acidic bile in a 30-foot cone. Each creature in the cone must make a DC 12 Dexterity saving throw, taking 40 (9d8) acid damage on a failed save, or half as much on a successful one. A creature reduced to 0 hit points by this damage is dissolved, leaving behind any objects it was carrying or wearing. Once used, the sprayer can't be used again.

A creature that interacts with the war machine spooks six invisible imps lurking in its bent chassis. The imps fly out and attack the source of their fright while yelling, "I can't go back to jail!" in Infernal. If the infernal war machine's acid sprayer hasn't already been used, one of the imps uses it.

A character who sits in the driver's seat notices a glowing Potion of Fire Resistance jammed in a dirty glove compartment along with six unpaid infernal parking tickets.

\subsection{E5: Planartarium Car}
\begin{DndReadAloud}{}
Glowing symbols on the floor line the perimeter of this domed car. In the middle of the car is a mechanical model suspended over a dais that has more glowing symbols on it. The model resembles a large wheel with a tall spire protruding from the top. A tiny brass ring floats above the spire's peak.

The spectral blue image of a modron materializes before the dais. It has a boxy shape and wears a pair of oversized glasses. The illusory modron adjusts its glasses and joyfully greets you in Common:

"Hello, and welcome to the Planartarium. You can learn more about the planes of existence by donating to this car. All proceeds benefit the establishment of order across the multiverse." Before disappearing, the modron gestures toward a short tube protruding from the floor near the dais.

\end{DndReadAloud}

The sixteen glowing sigils on the floor represent each of the Outer Planes, while the raised dais in the center of the room bears symbols of the four elemental Inner Planes. The wheel-shaped model above the dais represents the Outlands and its sixteen gate-towns. The tiny ring at the top of the spire represents Sigil.

A magical mote of golden light denotes the train's location. Since the train's travels are usually limited to the Inner Planes, the Outer Planes, and the Outlands, this mote can be seen floating above one of the symbols or above the model wheel representing the Outlands. If the train is in some other location not represented in this car, no mote is visible.

\subsubsection{Donation Tube}
The short tube rising from the floor near the dais has a slot large enough to insert coins, gemstones, and other small valuables. Such valuables are then teleported to the cart in area E9.

Whenever a creature drops at least 10 gp worth of treasure into the tube, Cosmo, the illusory modron, reappears. It thanks the donor and asks which plane of existence they'd like to learn about.

Cosmo can give a brief lecture on any \textit{plane of existence} described in the \textit{Dungeon Master's Guide}. If a character asks about a plane on the train's route, Cosmo also alludes to noteworthy planar effects. For example, if asked about the Nine Hells, the spectral duodrone might say, "Rest assured, the Concordant Express is reinforced to withstand extreme heat and extreme cold." Cosmo lectures about one plane per donation and vanishes after each lecture.

If a character inserts something worthless into the donation slot or otherwise tries to fool it with counterfeits and the like, the tube disgorges poisonous gas that quickly fills the compartment. This gas duplicates the effect of a \textit{stinking cloud} spell that lasts for 1 minute or until it is dispersed.

\subsection{E6: Passenger Car}
\begin{DndReadAloud}{}
A mind flayer in a beige coat kneels beside a body in the aisle of this passenger car. Three curious onlookers—a cambion, a drow, and bronze dwarf with flames for hair—poke their heads out of nearby cabins. "No one move a muscle!" warns the mind flayer. "This whole car just became a crime scene."

\end{DndReadAloud}

This car contains three cozy passenger cabins and two privies (one at each end of the car). Quintus Malvesh, an aasimar cartographer, lies dead in the aisle (in the square marked X on map 11.2), and a detective is already on the case.

\subsubsection{Illithid Detective}
The detective is a lawful neutral \textbf{mind flayer} named \textbf{Ignatius Inkblot}. Seeing the characters, Ignatius lights a pipe with a match it holds with one of its face tentacles. The pipe exudes a green vapor that smells of absinthe. Since the characters weren't on the car when the murder occurred, Ignatius doesn't regard them as suspects and is friendly toward them. Ignatius can speak Common but loathes doing so; the mind flayer prefers to converse telepathically. Its tone is strident, and it uses stereotypical detective phrases such as "The game is afoot!" and "The devil is in the details!"

Ignatius tries to enlist the characters in its investigation. If they agree to help, the mind flayer asks them to look for clues and question suspects. (In the meantime, the mind flayer quietly casts \textit{detect thoughts}, focusing on every creature in the Passenger Car.) Ignatius asks that no one leave the car until the mystery is solved. If the characters leave anyway, Ignatius solves the case without their help.

\subparagraph{What Ignatius Knows}
Ignatius has already determined that Quintus was killed while walking to his cabin from one of the Passenger Car's two privies. The aasimar's body lies outside the door to the middle cabin, which Quintus and another passenger shared. Ignatius was using the other privy at the time of Quintus's death and heard the aasimar gasp before hitting the floor. Ignatius quickly reached the scene, but the killer was nowhere in sight.

\subsubsection{Murder Suspects}
The suspects aboard the Passenger Car are as follows:

\textbf{Abernathy Vernus} (lawful evil \textbf{cambion}) is a friendly, ingratiating fellow who claims to have won a free tour aboard the Concordant Express. He insists that others call him Vern. No other passengers occupy his cabin (the one marked "a" on map 11.2), which Vern shares with a \textbf{monodrone} valet that answers to the name Higglesworth.

\textbf{Ethlynn Stalaczic} (neutral, drow \textbf{mage}) is a friendly, somewhat gullible dragonologist on a planar quest to find a "time dragon" (which may or may not be a real thing). Until a few moments ago, Ethlynn and Quintus were cabin mates. Ethlynn currently shares her cabin (the one marked "b" on map 11.2) with a \textbf{monodrone} valet called Bot.

\textbf{Meldar} (lawful neutral \textbf{azer}) is indifferent toward others, speaks nothing but Ignan, and provides curt answers to questions put to them. The azer is returning to the City of Brass after attending a smiths' convention in Bytopia. Meldar shares a cabin (the one marked "c" on map 11.2) with \textbf{Ignatius Inkblot} and a monodrone valet known as Orb.

The monodrones never leave the train. Ignatius was already aboard when Meldar boarded the train a day ago in Bytopia. The other suspects boarded the train near Excelsior (a gate-town in the Outlands) shortly before the characters came aboard.

\subsubsection{Murder Victim}
Characters can examine Quintus or cast spells on the corpse. The following information can be learned by doing so:


\begin{description}
\item[Examining the Corpse.]
A character who examines Quintus and succeeds on a DC 11 Intelligence (Arcana) or Wisdom (Medicine) check determines that the aasimar was killed by three bolts of magical force identical to those created by a \textit{magic missile} spell. Given where the body fell, the direction Quintus was walking when he died, and where the wounds on the body are located, a character can deduce, with a successful DC 16 Intelligence (Investigation) check, that the killer likely didn't come from Meldar's cabin.
\item[Raising the Dead.]
A character who has the \textit{raise dead} spell prepared can use it to restore Quintus to life. (The devas in areas E3 and E8 can also provide this service, if asked nicely.) Quintus is chaotic neutral and uses the \textbf{noble} stat block, although he is unarmored and carries no weapons. Alive, the aasimar is less cooperative than when he was dead, for reasons explained in the "Motive" section below. He offers no possible reasons for why someone would want to kill him.
\item[Speaking with the Dead.]
If a \textit{speak with dead} spell is cast on the corpse, the aasimar's spirit answers questions honestly. Quintus did not see his killer, but he knows one reason why every other passenger in this car (besides Ignatius) might want him dead. These reasons are summarized in the Murder Suspects table.
\end{description}

\subsubsection{Searching the Cabins}
Characters can search for clues in the passenger cabins and question the modron valet in each one:


\begin{description}
\item[Cabin A.]
This cabin holds Vern's battered suitcase, which contains two Traveler's Clothes, three odd-looking compasses, and a set of \textit{tinker's tools}. If asked about the compasses, Vern says he builds and repairs them to pass the time on long trips.
\item[Cabin B.]
Ethlynn shared this cabin with the murder victim and knows Quintus was a cartographer. Ethlynn confirms that Quintus left their cabin to use the privy at one end of the train car. Characters who search the cabin find Quintus's satchel, which contains two spare Traveler's Clothes, a set of \textit{cartographer's tools}, and a Map or Scroll Case containing what look like ancient treasure maps. A similar map, stolen by Ethlynn while Quintus was visiting the privy, is stuffed in her spidersilk handbag. (Ethlynn insists Quintus gave the map to her, which he never did. He did, however, tell her that the map shows the way to a time dragon's lair, which it does not.) A character who spends 1 minute examining any one of Quintus's maps can, with a successful DC 16 Intelligence (Investigation) check, cast serious doubt on the map's accuracy.
\item[Cabin C.]
Meldar and Ignatius share this cabin, which contains nothing incriminating. Its occupants travel light.
\item[Questioning the Modrons.]
Monodrone valets aboard the Concordant Express have built-in language interpreters instead of voice boxes; consequently, they can understand any spoken language but can't speak. Bot and Orb shrug their shoulders if questioned about the murder (of which they know nothing). Higglesworth simply points to itself if questioned about the murder.
\end{description}

\subsubsection{Solving the Mystery}
The monodrone Higglesworth killed Quintus using a \textit{wand of magic missiles}, which is hidden in a secret compartment inside the monodrone's spherical body. The monodrone won't subject itself to inspection and attacks anyone who tries to search or dismantle it. Its impeccable valet programming prevents it from incriminating passengers or accusing them of any wrongdoing.

Higglesworth did not commit the crime willingly. The cambion Vern used his tinker's tools and magical know-how to reprogram the monodrone valet to murder Quintus using a wand Vern pilfered from a wizard's apprentice in Excelsior. Vern then used the modron's body to hide the murder weapon and instructed Higglesworth not to tolerate invasive scrutiny.

\subparagraph{Murder Weapon}
A character can use an action to search Higglesworth for a secret compartment, finding it with a successful DC 13 Wisdom (Perception) check, but this can be done only while the monodrone is incapacitated. The same is true for opening the compartment, which requires another action, thieves' tools or tinker's tools, and a successful DC 11 Dexterity check. Opening the compartment causes the wand to tumble out.

If Higglesworth is reduced to 0 hit points, its body disintegrates, but the \textit{wand of magic missiles} hidden inside it remains.

One of the wand's charges has been expended—the same charge that killed Quintus.

\subparagraph{Motive}
Vern is a con artist who makes and sells "magic compasses" to gullible folk looking for an easy way to find that which their hearts desire. Quintus, also a con artist, made and sold bogus treasure maps. Vern and Quintus both identified Ethlynn as a likely mark, but Quintus got to her first by preying on her dream of finding a time dragon's lair. Vern decided to eliminate the competition while laying the blame on a malfunctioning modron that can't speak in its own defense.

\subparagraph{J'Accuse}
If the characters accuse Ethlynn or Meldar of murdering Quintus, Ignatius telepathically warns the characters that their investigation has gone "off the rails." He then recommends they "look for clues within clues" and "try to understand the motive behind the killing." Gifted with the ability to cast \textit{detect thoughts} at will, Ignatius has already identified Vern as the individual responsible for Quintus's death, but Ignatius wants the characters to see the truth for themselves—not simply take a mind flayer's word for it.

\subparagraph{Case Closed}
The mystery concludes with Ignatius confining Vern to his cabin until the authorities in Mechanus remove the cambion from the train. If the characters don't have the means to raise Quintus from the dead, Ignatius suggests they speak to the deva guarding the Jail Car and convince the angel to revive the dead aasimar. The characters are free to keep the \textit{wand of magic missiles}, since Ignatius doesn't need it to make a case against Vern.

Ignatius won't press Ethlynn on the theft of Quintus's bogus map, but characters may do so. Ethlynn can be guilted into returning the map and apologizing for the theft. As for Quintus, the aasimar won't admit to any wrongdoing unless magic is used to coerce him.

\begin{DndTable}[header=Murder Suspects]{lX}

Suspect & Possible Motive\\
Abernathy Vernus & The cambion seemed upset, maybe even jealous, that Ethlynn and Quintus were getting along so well. The cambion seemed to have his heart set on spending more time with Ethlynn, which necessitated getting rid of Quintus.\\
Ethlynn Stalaczic & Ethlynn was eyeballing Quintus's maps. Killing Quintus would allow her to steal the maps without anyone knowing.\\
Meldar & Quintus referred to Meldar as a "hothead" when he didn't think the azer was close enough to hear him. Killing Quintus could've been payback for the insult. (Since Meldar knows only one language—Ignan—and the insult wasn't spoken in that tongue, this motive evaporates quickly.)\\
\end{DndTable}

\subsection{E7: Temple Car}
\begin{DndReadAloud}{}
Stained-glass windows decorate the walls of this car, while its floors are adorned with exquisitely patterned rugs. Several columns support a vaulted mosaic ceiling. Kneeling before an altar is a serene priest in sunny robes.

\end{DndReadAloud}

The priest is, in fact, a \textbf{gray slaad} using its Change Shape ability to fool the characters. The slaad was sent here to kill them, plain and simple. It killed the real priest, stole his robes, and threw his body and holy symbol off the train.

The slaad works for an evil, multiverse-spanning organization called the Syndicate of Terror, Extortion, Assassination, and Larceny (STEAL), which has been spying on the characters for some time and wants their mission aboard the Concordant Express to fail. (This is doubly true if the characters are working for STEAL's rival, the Golden Vault.) If you're weaving multiple adventures from this anthology into your campaign, one or more villains from earlier adventures might have secret ties to this evil organization. Likely candidates (assuming they're still alive) include Quentin Togglepocket from "The Stygian Gambit," Guildmaster Dusk from "Masterpiece Imbroglio," and \textbf{Nixylanna Vidorant} from "Vidorant's Vault." If you're using the rivals described in this book's introduction, their patron might be a STEAL operative.

One of STEAL's operatives, having recently acquired the gray slaad's control gem, used the gem to command the slaad to slip aboard the Concordant Express and kill the characters on sight. The slaad loathes being surrounded by modrons on the train and fights to the death.

\subsubsection{Scrying Sensor}
The slaad's master watches the battle unfold through an invisible scrying sensor like the one created by a \textit{scrying} spell. If the characters defeat the slaad, a magically altered voice issues from the invisible sensor and says in Common, "So much for that. Next time you won't be so lucky!" If the characters are working for the Golden Vault, the voice adds, "The Golden Vault won't be able to protect you from our wrath." The scrying sensor then disappears.

\subsection{E8: Jail Car}
This soundproof car is two stories tall (30 feet high) and encased in thick iron. Four pentadrones guard the rooftop. These guards are positioned near the corners of the roof, allowing them to spot trouble coming from the sides of the car.

The outer doors of the car are locked tight, as is the 10-foot-square skylight in the roof. As an action, a character with thieves' tools can try to open one of the doors or the skylight, doing so with a successful DC 17 Dexterity check. The skylight can also be shattered; it is made of thick glass and has AC 13, 27 hit points, and immunity to poison and psychic damage. Conversely, the doors are too strong to be forced open by anything other than \textit{knock} spells or similar magic.

When the characters enter the car, read or paraphrase the following:

\begin{DndReadAloud}{}
Two iron balconies protrude from the side walls of this train car. Iron stairwells connect these balconies to the lower level. Lining the walls on both levels are cell doors made of riveted iron, each one fitted with a large, gear-shaped keyhole. Pale light shines through paper-thin gaps around the doors, but not the keyholes.

A stocky angel with cerulean skin and a pointed helmet leans against the railing on one balcony, a radiant flanged mace at its hip and its fiery wings folded behind it.

\end{DndReadAloud}

\subsubsection{Omid the Deva}
The angel guarding the Jail Car is Omid, a \textbf{deva} with fiery instead of feathery wings. Omid has orders to deliver \textbf{the Stranger} safely to Mechanus, where judgment awaits.

The deva's glowing mace is reminiscent of a large key. The head of the mace lines up neatly with the gear-shaped keyhole on each of the car's eight cell doors. As an action, a character can use the mace-key to unlock any one of these doors, which can be opened no other way.

Omid remains in the Jail Car and can't be lured from it. When the deva encounters the characters, read the following:

\begin{DndReadAloud}{}
The angel unfurls its wings, releasing a flurry of glowing embers, and says, "Speak your business, trespassers, and speak it truthfully."

\end{DndReadAloud}

The characters need not resort to violence to obtain the deva's mace-key. Here are some alternate approaches:


\begin{description}
\item[Clever Deception.]
The characters can pose as emissaries of a jurisdiction with greater claim to \textbf{the Stranger}'s custody. Convincing Omid to release \textbf{the Stranger} requires a successful DC 17 group Charisma (Deception) check. If the characters used a forgery kit to create documentation to support their claim, the group check is made with advantage.
\item[Convincing Argument.]
Since Omid knows little about \textbf{the Stranger}, a character can convince the deva that the outlaw has been wrongly accused and should be released. Making the argument might take a minute or two, after which the character must succeed on a DC 17 Charisma (Persuasion) check to secure \textbf{the Stranger}'s release.
\item[Friendly Game.]
A character who has a gaming set can challenge the deva to a friendly game. Omid loves games and is too proud (and bored) not to accept. If the character wins the game, Omid allows the party to speak with \textbf{the Stranger} (but not set them free). Determine the game's outcome by having the deva's challenger make a DC 19 Intelligence check. Proficiency with the gaming set allows the character to add their proficiency bonus to their roll. If the check is successful, the character wins. On a failed check, the character can opt to cheat in the late stages of the game and make a DC 19 Charisma (Deception) check. A successful check overrides the earlier failure. If the character loses gracefully and didn't cheat, the deva (a gracious winner) still allows the party to speak with \textbf{the Stranger}.
\end{description}

\subsubsection{Cells}
The door to each cell opens into a 10-foot-square extradimensional space brightly lit by magic. The cells and their contents are as follows:


\begin{description}
\item[Cell 1.]
This cell contains a hostile \textbf{nycaloth} known as Dardo (not its true name). If someone opens the door to its cell, Dardo tries to exit the cell on its next turn, attacking anyone who stands in its way. The nycaloth then tries to leave the train. \textbf{The Stranger} (see below) happens to know Dardo's true name: a simple math equation. The characters can use this knowledge to keep the nycaloth from attacking them, or to force it to work for them.
\item[Cell 2.]
This cell is empty.
\item[Cell 3.]
A bench against the far wall has a \textit{+2 rod of the pact keeper} resting on it. The cell is otherwise empty.
\item[Cell 4.]
This cell contains a friendly \textbf{erinyes} who recently helped an efreeti named Vrakir get his hands on the Book of Vile Darkness (Variant). Speaking with this prisoner might kick off "Fire and Darkness," another adventure in this book. \textbf{The Stranger} (see below) doesn't know this erinyes's true name.
\item[Cell 5.]
This cell is empty.
\item[Cell 6.]
This cell is empty.
\item[Cell 7.]
A \textbf{monodrone} janitor accidentally locked itself in this cell.
\item[Cell 8.]
This cell contains \textbf{the Stranger} (see below).
\end{description}

\subsubsection{The Stranger}
This individual looks like a human with a flat-brimmed leather hat and a tattered traveler's cloak. \textbf{The Stranger}'s face is permanently obscured behind a shifting blur like that of a deep desert haze. \textbf{The Stranger} is chaotic neutral and uses the \textbf{archmage} stat block, though they abhor combat. \textbf{The Stranger} wears an invisible \textit{ring of mind shielding} that thwarts magical attempts to read their thoughts. \textbf{The Stranger} also wears \textit{dimensional shackles} but is eager to be rid of them. One touch from Omid's mace-key causes the shackles to fall off, freeing \textbf{the Stranger}.

\textbf{The Stranger} is initially indifferent to the characters but becomes friendly if presented with the prospect of escape, at which point \textbf{the Stranger} replies in a gravelly voice:

\begin{DndReadAloud}{}
"Let me guess. You want a name. Maybe more than one. Got a pen?"

\end{DndReadAloud}

\textbf{The Stranger} promises to recite the true names the characters seek in exchange for freedom but warns them of the heavy toll such knowledge carries.

If the characters are after the true names of Karnyros, Errtok, and Hexalanthe, \textbf{the Stranger} reveals the following once free of the \textit{dimensional shackles} and outside the cell:


\begin{itemize}
\item Karnyros's true name is a series of guttural growls and grimaces, followed by a sound like a popping balloon.
\item Errtok's true name is Ar-lothe Gothu Ka.
\item Hexalanthe's true name is a low-toned, rhythmic chant accompanied by a series of hand gestures, dance moves, and drumrolls.
\end{itemize}

If the characters are after the true name of some other creature, perhaps a fiendish nemesis or tenuous planar ally, \textbf{the Stranger} might know that creature's true name as well and can reveal it to the characters. Even if \textbf{the Stranger} doesn't know a particular true name, they can give characters the true names of one or more Fiends who have the power to help the characters in some other way.

If you have trouble coming up with ideas for new true names, roll on the Random True Names table at the end of this adventure when you need one.

\subsection{E9: Engine Car}
\begin{DndReadAloud}{}
A massive furnace dominates the interior of this noisy car. Two spherical modrons shovel coins into the roaring fire while a brawny, nine-foot-tall automaton guards a treasure-filled fuel cart nearby. "Faster!" bellows a deep, disembodied voice.

\end{DndReadAloud}

The sapient Engine Car is the source of the deep, booming voice. The car speaks Common and Modron. While it's moving and not talking, the Engine Car labors through an endless chain of short, forceful breaths timed with each rhythmic churn of its mechanisms. Its undercarriage chatters as dozens of mechanical arms lay railway tracks before the train.

Inside the car, two monodrones stoke the fire in the furnace by shoveling loot into it from a nearby cart (see "Treasure" below). A \textbf{shield guardian} stands next to the cart, guarding the treasure from thieves. Thanks to its Spell Storing trait, the guardian can cast \textit{sleep} (4th-level version) once.

\subsubsection{Furnace}
Any object not being worn or carried that is placed inside the blazing furnace is incinerated instantly. Even magic items succumb to the blazing furnace, though the flames aren't powerful enough to harm artifacts.

If the furnace isn't replenished every hour with at least 100 gp worth of treasure, the Engine Car sputters and the train comes to a halt, causing a guard patrol (see "Guard Patrols" earlier in the adventure) to investigate.

A creature that enters the furnace for the first time on a turn or starts its turn there takes 55 (10d10) fire damage. A creature reduced to 0 hit points by this damage dies and is reduced to a pile of ash.

The furnace door can be latched from the outside. A creature trapped inside the furnace can use an action to try to force open the door, doing so with a successful DC 20 Strength (Athletics) check.

\subsubsection{Treasure}
The treasure cart contains 3,500 gp, 6,000 sp, and 1,200 cp. Scattered among the coins are the following items:


\begin{itemize}
\item Five Amethyst (100 gp each)
\item Gilded mind flayer skull (150 gp)
\item Horn of Valhalla, Silver
\item Suit of Mithral Chain Mail
\item Potion of Greater Healing
\item Silvered \textit{+3 greatsword} that once belonged to a githyanki knight
\end{itemize}

\section{Conclusion}
After imparting the true names, the freed Stranger tips their hat and leaps from the train. Once out of range of the train's teleportation ward, \textbf{the Stranger} casts \textit{plane shift} and disappears. The characters can leave the Concordant Express in a similar fashion, or they can remain on the train until it arrives in Mechanus and secure safe passage home from there. Whether or not they helped solve the murder in the Passenger Car (area E7), characters who disembark in Mechanus see \textbf{Ignatius Inkblot}, the mind flayer detective, escorting the cambion Abernathy Vernus off the train and into the waiting arms of a group of modrons.

\subsection{For the Golden Vault}
If the characters are working for the Golden Vault, the organization sends the group's contact to retrieve the true names the day after the characters record them. This operative offers the characters a rare magic item of their choice (subject to your approval) as payment, which is delivered to the characters the next day.

\subsection{Other Rewards}
If the characters helped solve the murder mystery aboard the Concordant Express, the mind flayer detective \textbf{Ignatius Inkblot} uses its \textit{plane shift} spell to deliver a wrapped parcel to the characters several days later. The parcel contains a token of the illithid's esteem: a masterfully crafted magnifying glass worth 1,000 gp.

Characters can also leverage any true names they learn into further rewards, as discussed below:


\begin{description}
\item[Home Ownership.]
Knowing a creature's true name may entitle the characters to that creature's property, such as a parcel of land, lair, or stronghold.
\item[Slayer's Weapon.]
A master smith might be able to work the true name of a creature into an \textit{arrow of slaying} or other target-specific magic item.
\end{description}

\subsection{Further Adventures}
The characters might have made a number of planar enemies during this heist. Subsequent adventures could include the following:

\textbf{STEAL}. The rival organization STEAL (see area E7) might figure prominently in future adventures. At higher levels, run-ins with STEAL operatives or raids on STEAL vaults might place the characters in conflict with the syndicate's leaders, including arcanaloths and ultroloths with true names only \textbf{the Stranger} knows.


\begin{description}
\item[Trials and Tribulations.]
The characters might be forced to stand before a tribunal in Mechanus for aiding an interplanar fugitive or answer to a powerful Celestial for slaying an angel.
\item[True Name Terror.]
Characters who learn the true names of powerful Fiends are marked for death by those Fiends, which send powerful underlings to slay the characters.
\end{description}

\begin{DndTable}[header=Random True Names]{cX}

d100 & True Name\\
01–03 & Sequence of musical notes played on a lyre\\
04–05 & "Zow Chokk Garthagua"\\
06–07 & "Lureff Zadil"\\
08–09 & Sound of a bucket dropped on the floor\\
10–11 & "Batridd Morath"\\
12–13 & "Agh Gizzirah Korondis"\\
14–15 & Third line of an obscure Auran lullaby\\
16–17 & "Naelurec" ("cerulean" backward)\\
18–19 & "Si Uulth Uurek"\\
20–21 & Sound of four eggs sizzling in a frying pan\\
22–23 & "Lamoura"\\
24–25 & "Haraknar Orvorag"\\
26–27 & Nine-second hum of a tuning fork\\
28–29 & "Kyrvaash Cesvyr"\\
30–31 & "Tsul Vandar"\\
32–33 & Sound of three pages tearing from a book\\
34–35 & "Vrah Koruun"\\
36–37 & "Jal-Shoth"\\
38–39 & Seven tolls of an iron bell\\
40–41 & "Barak-Shivad of the Night Realm"\\
42–43 & "Winterglory"\\
44–45 & Short whistle, followed by two long ones\\
46–47 & "Kelefis the Thrice-Hated"\\
48–49 & "Nefarion"\\
50–52 & Three specific prime numbers strung together\\
53–54 & "Zizz-Ganneth"\\
55–56 & "Red Rust"\\
57–59 & Sequence of red and blue lights flashing\\
60–61 & "Xurgranach"\\
62–63 & "Enchridus 3291"\\
64–65 & Last spoken words of a long-dead hero\\
66–67 & "Valmorag the White"\\
68–69 & "Suzakiro Arzuun of the Rusty Blades"\\
70–71 & Sound of fingers drumming on a table\\
72–73 & "Myranna Vos Aeldar"\\
74–75 & "Kzokzys"\\
76–77 & Gasp of someone stabbed through the heart\\
78–79 & "U'chud, the Eye of Zefir-Zaskos"\\
80–81 & "Quas Quannok"\\
82–83 & Sad mewling of a kitten\\
84–85 & "Skuzrayle"\\
86–87 & "Nebek Velflam Von Qwizzid"\\
88–89 & Sound of a lute being smashed to pieces\\
90–91 & "Old Sinisyphere"\\
92–93 & "Ghax'ru the Abominable"\\
94–96 & Sound of three acorns crunched between teeth\\
97–98 & "Knob"\\
99–00 & "The Nameless One"\\
\end{DndTable}

\part{Magic Items}
\DndItemHeader{+2 Rod of the Pact Keeper}{Rare (requires attunement by a warlock)}
While holding this rod, you gain a +2 bonus to spell attack rolls and to the saving throw DCs of your warlock spells.

In addition, you can regain one warlock spell slot as an action while holding the rod. You can't use this property again until you finish a long rest.

\DndItemHeader{Amethyst}{}
A deep purple gemstone.

\DndItemHeader{Book of Vile Darkness}{Wondrous item, Artifact (requires attunement)}
The contents of this foul manuscript of ineffable wickedness are the meat and drink of those in evil's thrall. No mortal was meant to know the secrets it contains, knowledge so horrid that to even glimpse the scrawled pages invites madness.

Most believe the lich-god Vecna authored the Book of Vile Darkness. He recorded in its pages every diseased idea, every unhinged thought, and every example of blackest magic he came across or devised. Vecna covered every vile topic he could, making the book a gruesome catalog of all mortal wrongs.

Other practitioners of evil have held the book and added their own input to its catalog of vile knowledge. Their additions are clear, for the writers of later works stitched whatever they were writing into the tome or, in some cases, made notations and additions to existing text. There are places where pages are missing, torn. or covered so completely with ink, blood, and scratches that the original text can't be divined.

Nature can't abide the book's presence. Ordinary plants wither in its presence, animals are unwilling to approach it, and the book gradually destroys whatever it touches. Even stone cracks and turns to powder if the book rests on it long enough.

A creature attuned to the book must spend 80 hours reading and studying it to digest its contents and reap its benefits. The creature can then freely modify the book's contents, provided that those modifications advance evil and expand the lore already contained within.

Whenever a non-evil creature attunes to the Book of Vile Darkness, that creature must make a DC 17 Charisma saving throw. On a failed save, the creature's alignment changes to neutral evil.

The Book of Vile Darkness remains with you only as long as you strive to work evil in the world. If you fail to perform at least one evil act within the span of 10 days, or if you willingly perform a good act, the book disappears. If you die while attuned to the book, an entity of great evil claims your soul. You can't be restored to life by any means while your soul remains imprisoned.

\subparagraph{Random Properties}
The Book of Vile Darkness has the following random properties:


\begin{itemize}
\item 3 Artifact Properties; Minor Beneficial Properties
\item 1 Artifact Properties; Major Beneficial Properties
\item 3 Artifact Properties; Minor Detrimental Properties
\item 2 Artifact Properties; Major Detrimental Properties
\end{itemize}

\subparagraph{Adjusted Ability Scores}
After you spend the requisite amount of time reading and studying the book, one ability score of your choice increases by 2, to a maximum of 24. Another ability score of your choice decreases by 2, to a minimum of 3. The book can't adjust your ability scores again.

\subparagraph{Mark of Darkness}
After you spend the requisite amount of time reading and studying the book, you acquire a physical disfigurement as a hideous sign of your devotion to vile darkness. An evil rune might appear on your face, your eyes might become glossy black, or horns might sprout from your forehead. Or you might become wizened and hideous, lose all facial features, gain a forked tongue, or some other feature the DM chooses. The mark of darkness grants you advantage on Charisma (Persuasion) checks made to interact with evil creatures and Charisma (Intimidation) checks made to interact with non-evil creatures.

\subparagraph{Command Evil}
While you are attuned to the book and holding it, you can use an action to cast the \textit{dominate monster} spell on an evil target (save DC 18). You can't use this property again until the next dawn.

\subparagraph{Dark Lore}
You can reference the Book of Vile Darkness whenever you make an Intelligence check to recall information about some aspect of evil, such as lore about demons. When you do so, double your proficiency bonus on that check.

\subparagraph{Dark Speech}
While you carry the Book of Vile Darkness and are attuned to it, you can use an action to recite words from its pages in a foul language known as Dark Speech. Each time you do so, you take 1d12 psychic damage, and each non-evil creature within 15 feet of you takes 3d6 psychic damage.

\subparagraph{Destroying the Book}
The Book of Vile Darkness allows pages to be torn from it, but any evil lore contained on those pages finds its way back into the book eventually, usually when a new author adds pages to the tome.

If a \textbf{solar} tears the book in two, the book is destroyed for 1d100 years, after which it reforms in some dark corner of the multiverse.

A creature attuned to the book for one hundred years can unearth a phrase hidden in the original text that, when translated to Celestial and spoken aloud, destroys both the speaker and the book in a blinding flash of radiance. However, as long as evil exists in the multiverse, the book reforms 1d10 × 100 years later.

If all evil in the multiverse is wiped out, the book turns to dust and is forever destroyed.

\DndItemHeader{Horn of Valhalla, Silver}{Wondrous item, Rare}
You can take a Magic action to blow this horn. In response, warrior spirits from the plane of Ysgard appear in unoccupied spaces within 60 feet of you. Each spirit uses the \textbf{Berserker} stat block and returns to Ysgard after 1 hour or when it drops to 0 Hit Points. The spirits look like living, breathing warriors, and they have Immunity to the Charmed and Frightened conditions. Once you use the horn, it can't be used again until 7 days have passed.

A silver horn summons 2 Berserkers. They are Friendly [Attitude] to you and your allies and follow your commands.

\DndItemHeader{Map or Scroll Case}{}
This cylindrical leather case can hold up to ten rolled-up paper (one sheet) or five rolled-up parchment (one sheet).

\DndItemHeader{Potion of Fire Resistance}{Uncommon}
When you drink this potion, you have Resistance to {{getFullImmRes item.resist}} damage for 1 hour.

\DndItemHeader{Potion of Greater Healing}{Uncommon}
You regain 4d4 + 4 Hit Points when you drink this potion. The potion's red liquid glimmers when agitated.

\DndItemHeader{Amulet of the Planes}{Wondrous item, Very Rare (requires attunement)}
While wearing this amulet, you can use an action to name a location that you are familiar with on another plane of existence. Then make a DC 15 Intelligence check. On a successful check, you cast the \textit{plane shift} spell. On a failure, you and each creature and object within 15 feet of you travel to a random destination. Roll a d100. On a 1-60, you travel to a random location on the plane you named. On a 61-100, you travel to a randomly determined plane of existence.

\DndItemHeader{Decanter of Endless Water}{Wondrous item, Uncommon}
This stoppered flask sloshes when shaken, as if it contains water. The decanter weighs 2 pounds.

You can use an action to remove the stopper and speak one of three command words, whereupon an amount of fresh water or salt water (your choice) pours out of the flask. The water stops pouring out at the start of your next turn. Choose from the following options:


\begin{itemize}
\item "Stream" produces 1 gallon of water.
\item "Fountain" produces 5 gallons of water.
\item "Geyser" produces 30 gallons of water that gushes forth in a geyser 30 feet long and 1 foot wide. As a bonus action while holding the decanter, you can aim the geyser at a creature you can see within 30 feet of you. The target must succeed on a DC 13 Strength saving throw or take 1d4 bludgeoning damage and fall prone. Instead of a creature, you can target an object that isn't being worn or carried and that weighs no more than 200 pounds. The object is either knocked over or pushed up to 15 feet away from you.
\end{itemize}

\DndItemHeader{Dimensional Shackles}{Wondrous item, Rare}
You can use an action to place these shackles on an incapacitated creature. The shackles adjust to fit a creature of Small to Large size. In addition to serving as mundane manacles, the shackles prevent a creature bound by them from using any method of extradimensional movement, including teleportation or travel to a different plane of existence. They don't prevent the creature from passing-through an interdimensional portal.

You and any creature you designate when you use the shackles can use an action to remove them. Once every 30 days, the bound creature can make a DC 30 Strength (Athletics) check. On a success, the creature breaks free and destroys the shackles.

\DndItemHeader{Ring of Mind Shielding}{Uncommon (requires attunement)}
While wearing this ring, you are immune to magic that allows other creatures to read your thoughts, determine whether you are lying, know your alignment, or know your creature type. Creatures can telepathically communicate with you only if you allow it.

You can use an action to cause the ring to become invisible until you use another action to make it visible, until you remove the ring, or until you die.

If you die while wearing the ring, your soul enters it, unless it already houses a soul. You can remain in the ring or depart for the afterlife. As long as your soul is in the ring, you can telepathically communicate with any creature wearing it. A wearer can't prevent this telepathic communication.

\DndItemHeader{Wand of Magic Missiles}{Uncommon}
This wand has 7 charges. While holding it, you can use an action to expend 1 or more of its charges to cast the \textit{magic missile} spell from it. For 1 charge, you cast the 1st-level version of the spell. You can increase the spell slot level by one for each additional charge you expend.

The wand regains 1d6 + 1 expended charges daily at dawn. If you expend the wand's last charge, roll a d20. On a 1, the wand crumbles into ashes and is destroyed.

\part*{Printable Statblocks}
\setlist[description]{nolistsep,listparindent=3pt,topsep=6pt,font={\bfseries\sffamily\small}}
\pagestyle{empty}

\begin{DndMonster}[float=t]{Archmage}
\DndMonsterType{Medium humanoid (any race), A}
\DndMonsterBasics[
armor-class = {12 (15 with \textit{mage armor})},
hit-points = {\DndDice{18d8 + 18}},
speed = {30 ft.}
]
\DndMonsterAbilityScores[
 str = 10,
 dex = 14,
 con = 12,
 int = 20,
 wis = 15,
 cha = 16
]
\DndMonsterDetails[
saving-throws = {Int +9, Wis +6},
skills = {Arcana +13, History +13},
damage-resistances = {damage from spells; bludgeoning, piercing, slashing (from stoneskin)},
languages = {any six languages},
challenge = 12,
]
\DndMonsterAction{Magic Resistance}
The archmage has advantage on saving throws against spells and other magical effects.

\DndMonsterAction{Spellcasting}
The archmage is an 18th-level spellcaster. Its spellcasting ability is Intelligence (spell save DC 17, 9 to hit with spell attacks). The archmage can cast \textit{disguise self} and \textit{invisibility} at will and has the following wizard spells prepared:
\begin{DndMonsterSpells}
  \DndInnateSpellLevel{disguise self, \textit{invisibility}}
  \DndMonsterSpellLevel{fire bolt, \textit{light}, \textit{mage hand}, \textit{prestidigitation}, \textit{shocking grasp}}
   \DndMonsterSpellLevel[1][4]{detect magic, \textit{identify}, \textit{mage armor}*, \textit{magic missile}}
   \DndMonsterSpellLevel[2][3]{detect thoughts, \textit{mirror image}, \textit{misty step}}
   \DndMonsterSpellLevel[3][3]{counterspell, \textit{fly}, \textit{lightning bolt}}
   \DndMonsterSpellLevel[4][3]{banishment, \textit{fire shield}, \textit{stoneskin}*}
   \DndMonsterSpellLevel[5][3]{cone of cold, \textit{scrying}, \textit{wall of force}}
   \DndMonsterSpellLevel[6][1]{globe of invulnerability}
   \DndMonsterSpellLevel[7][1]{teleport}
   \DndMonsterSpellLevel[8][1]{mind blank*}
\end{DndMonsterSpells}
\DndMonsterSection{Actions}
\DndMonsterAction{Dagger}
\textsl{Melee or Ranged Weapon Attack:} 6 to hit, reach 5 ft. or range 20/60 ft., one target. \textsl{Hit:} 4 (1d4 + 2) piercing damage.

\end{DndMonster}



\begin{DndMonster}[float=t]{Azer}
\DndMonsterType{Medium elemental, lawful neutral}
\DndMonsterBasics[
armor-class = {17 (natural armor)},
hit-points = {\DndDice{6d8 + 12}},
speed = {30 ft.}
]
\DndMonsterAbilityScores[
 str = 17,
 dex = 12,
 con = 15,
 int = 12,
 wis = 13,
 cha = 10
]
\DndMonsterDetails[
saving-throws = {Con +4},
damage-immunities = {fire, poison},
condition-immunities = {poisoned},
languages = {Ignan},
challenge = 2,
]
\DndMonsterAction{Heated Body}
A creature that touches the azer or hits it with a melee attack while within 5 feet of it takes 5 (1d10) fire damage.

\DndMonsterAction{Heated Weapons}
When the azer hits with a metal melee weapon, it deals an extra 3 (1d6) fire damage (included in the attack).

\DndMonsterAction{Illumination}
The azer sheds bright light in a 10-foot radius and dim light for an additional 10 feet.

\DndMonsterSection{Actions}
\DndMonsterAction{Warhammer}
\textsl{Melee Weapon Attack:} 5 to hit, reach 5 ft., one target. \textsl{Hit:} 7 (1d8 + 3) bludgeoning damage, or 8 (1d10 + 3) bludgeoning damage if used with two hands to make a melee attack, plus 3 (1d6) fire damage.

\end{DndMonster}



\begin{DndMonster}[float=t]{Berserker}
\DndMonsterType{Small or Medium humanoid, A}
\DndMonsterBasics[
armor-class = {13},
hit-points = {\DndDice{9d8 + 27}},
speed = {30 ft.}
]
\DndMonsterAbilityScores[
 str = 16,
 dex = 12,
 con = 17,
 int = 9,
 wis = 11,
 cha = 9
]
\DndMonsterDetails[
languages = {Common},
challenge = 2,
]
\DndMonsterSection{Actions}
\DndMonsterAction{Greataxe}
m +5 (with Advantage if the target doesn't have all its Hit Points), reach 5 ft. \textsl{Hit:} 9 (1d12 + 3) Slashing damage.

\end{DndMonster}



\begin{DndMonster}[float=t]{Cambion}
\DndMonsterType{Medium fiend, any evil alignment}
\DndMonsterBasics[
armor-class = {19 (\textit{scale mail})},
hit-points = {\DndDice{11d8 + 33}},
speed = {30 ft., fly 60 ft.}
]
\DndMonsterAbilityScores[
 str = 18,
 dex = 18,
 con = 16,
 int = 14,
 wis = 12,
 cha = 16
]
\DndMonsterDetails[
saving-throws = {Str +7, Con +6, Int +5, Cha +6},
skills = {Deception +6, Intimidation +6, Perception +4, Stealth +7},
damage-resistances = {cold; fire; lightning; poison; bludgeoning, piercing, slashing from nonmagical attacks},
senses = {darkvision 60 ft., passive perception 14},
languages = {Abyssal, Common, Infernal},
challenge = 5,
]
\DndMonsterAction{Fiendish Blessing}
The AC of the cambion includes its Charisma bonus.

\DndMonsterAction{Innate Spellcasting}
The cambion's spellcasting ability is Charisma (spell save DC 14). The cambion can innately cast the following spells, requiring no material components:
\begin{DndMonsterSpells}
  \DndInnateSpellLevel[]{plane shift (self only)}
  \DndInnateSpellLevel[3]{alter self, \textit{command}, \textit{detect magic}}
\end{DndMonsterSpells}
\DndMonsterSection{Actions}
\DndMonsterAction{Multiattack}
The cambion makes two melee attacks or uses its Fire Ray twice.

\DndMonsterAction{Spear}
\textsl{Melee or Ranged Weapon Attack:} 7 to hit, reach 5 ft. or range 20/60 ft., one target. \textsl{Hit:} 7 (1d6 + 4) piercing damage, or 8 (1d8 + 4) piercing damage if used with two hands to make a melee attack, plus 3 (1d6) fire damage.

\DndMonsterAction{Fire Ray}
\textsl{Ranged Spell Attack:} 7 to hit, range 120 ft., one target. \textsl{Hit:} 10 (3d6) fire damage.

\DndMonsterAction{Fiendish Charm}
One humanoid the cambion can see within 30 feet of it must succeed on a DC 14 Wisdom saving throw or be magically charmed for 1 day. The charmed target obeys the cambion's spoken commands. If the target suffers any harm from the cambion or another creature or receives a suicidal command from the cambion, the target can repeat the saving throw, ending the effect on itself on a success. If a target's saving throw is successful, or if the effect ends for it, the creature is immune to the cambion's Fiendish Charm for the next 24 hours.

\end{DndMonster}



\begin{DndMonster}[float=t]{Deva}
\DndMonsterType{Medium celestial, lawful good}
\DndMonsterBasics[
armor-class = {17 (natural armor)},
hit-points = {\DndDice{16d8 + 64}},
speed = {30 ft., fly 90 ft.}
]
\DndMonsterAbilityScores[
 str = 18,
 dex = 18,
 con = 18,
 int = 17,
 wis = 20,
 cha = 20
]
\DndMonsterDetails[
saving-throws = {Wis +9, Cha +9},
skills = {Insight +9, Perception +9},
damage-resistances = {radiant; bludgeoning, piercing, slashing from nonmagical attacks},
condition-immunities = {charmed, exhaustion, frightened},
senses = {darkvision 120 ft., passive perception 19},
languages = {all, telepathy 120 ft.},
challenge = 10,
]
\DndMonsterAction{Angelic Weapons}
The deva's weapon attacks are magical. When the deva hits with any weapon, the weapon deals an extra 4d8 radiant damage (included in the attack).

\DndMonsterAction{Magic Resistance}
The deva has advantage on saving throws against spells and other magical effects.

\DndMonsterAction{Innate Spellcasting}
The deva's spellcasting ability is Charisma (spell save DC 17). The deva can innately cast the following spells, requiring only verbal components:
\begin{DndMonsterSpells}
  \DndInnateSpellLevel{detect evil and good}
  \DndInnateSpellLevel[1]{commune, \textit{raise dead}}
\end{DndMonsterSpells}
\DndMonsterSection{Actions}
\DndMonsterAction{Multiattack}
The deva makes two melee attacks.

\DndMonsterAction{Mace}
\textsl{Melee Weapon Attack:} 8 to hit, reach 5 ft., one target. \textsl{Hit:} 7 (1d6 + 4) bludgeoning damage plus 18 (4d8) radiant damage.

\DndMonsterAction{Healing Touch (3/Day)}
The deva touches another creature. The target magically regains 20 (4d8 + 2) hit points and is freed from any curse, disease, poison, blindness, or deafness.

\DndMonsterAction{Change Shape}
The deva magically polymorphs into a humanoid or beast that has a challenge rating equal to or less than its own, or back into its true form. It reverts to its true form if it dies. Any equipment it is wearing or carrying is absorbed or borne by the new form (the deva's choice).

In a new form, the deva retains its game statistics and ability to speak, but its AC, movement modes, Strength, Dexterity, and special senses are replaced by those of the new form, and it gains any statistics and capabilities (except class features, legendary actions, and lair actions) that the new form has but that it lacks.

\end{DndMonster}



\begin{DndMonster}[float=t]{Duodrone}
\DndMonsterType{Medium construct, lawful neutral}
\DndMonsterBasics[
armor-class = {15 (natural armor)},
hit-points = {\DndDice{2d8 + 2}},
speed = {30 ft.}
]
\DndMonsterAbilityScores[
 str = 11,
 dex = 13,
 con = 12,
 int = 6,
 wis = 10,
 cha = 7
]
\DndMonsterDetails[
senses = {truesight 120 ft., passive perception 10},
languages = {Modron},
challenge = 1/4,
]
\DndMonsterAction{Axiomatic Mind}
The duodrone can't be compelled to act in a manner contrary to its nature or its instructions.

\DndMonsterAction{Disintegration}
If the duodrone dies, its body disintegrates into dust, leaving behind its weapons and anything else it was carrying.

\DndMonsterSection{Actions}
\DndMonsterAction{Multiattack}
The duodrone makes two fist attacks or two javelin attacks.

\DndMonsterAction{Fist}
\textsl{Melee Weapon Attack:} 2 to hit, reach 5 ft., one target. \textsl{Hit:} 2 (1d4) bludgeoning damage.

\DndMonsterAction{Javelin}
\textsl{Melee or Ranged Weapon Attack:} 3 to hit, reach 5 ft. or range 30/120 ft., one target. \textsl{Hit:} 4 (1d6 + 1) piercing damage.

\end{DndMonster}



\begin{DndMonster}[float=t]{Erinyes}
\DndMonsterType{Medium fiend (devil), lawful evil}
\DndMonsterBasics[
armor-class = {18 (\textit{plate armor})},
hit-points = {\DndDice{18d8 + 72}},
speed = {30 ft., fly 60 ft.}
]
\DndMonsterAbilityScores[
 str = 18,
 dex = 16,
 con = 18,
 int = 14,
 wis = 14,
 cha = 18
]
\DndMonsterDetails[
saving-throws = {Dex +7, Con +8, Wis +6, Cha +8},
damage-resistances = {cold; bludgeoning, piercing, slashing from nonmagical attacks that aren't silvered},
damage-immunities = {fire, poison},
condition-immunities = {poisoned},
senses = {truesight 120 ft., passive perception 12},
languages = {Infernal, telepathy 120 ft.},
challenge = 12,
]
\DndMonsterAction{Hellish Weapons}
The erinyes's weapon attacks are magical and deal an extra 13 (3d8) poison damage on a hit (included in the attacks).

\DndMonsterAction{Magic Resistance}
The erinyes has advantage on saving throws against spells and other magical effects.

\DndMonsterSection{Actions}
\DndMonsterAction{Multiattack}
The erinyes makes three attacks.

\DndMonsterAction{Longsword}
\textsl{Melee Weapon Attack:} 8 to hit, reach 5 ft., one target. \textsl{Hit:} 8 (1d8 + 4) slashing damage, or 9 (1d10 + 4) slashing damage if used with two hands, plus 13 (3d8) poison damage.

\DndMonsterAction{Longbow}
\textsl{Ranged Weapon Attack:} 7 to hit, range 150/600 ft., one target. \textsl{Hit:} 7 (1d8 + 3) piercing damage plus 13 (3d8) poison damage, and the target must succeed on a DC 14 Constitution saving throw or be poisoned. The poison lasts until it is removed by the \textit{lesser restoration} spell or similar magic.

\DndMonsterSection{Reactions}
\DndMonsterAction{Parry}
The erinyes adds 4 to its AC against one melee attack that would hit it. To do so, the erinyes must see the attacker and be wielding a melee weapon.

\end{DndMonster}



\begin{DndMonster}[float=t]{Gray Slaad}
\DndMonsterType{Medium aberration (shapechanger), chaotic neutral}
\DndMonsterBasics[
armor-class = {18 (natural armor)},
hit-points = {\DndDice{17d8 + 51}},
speed = {30 ft.}
]
\DndMonsterAbilityScores[
 str = 17,
 dex = 17,
 con = 16,
 int = 13,
 wis = 8,
 cha = 14
]
\DndMonsterDetails[
skills = {Arcana +5, Perception +7},
damage-resistances = {acid, cold, fire, lightning, thunder},
senses = {blindsight 60 ft., darkvision 60 ft., passive perception 17},
languages = {Slaad, telepathy 60 ft.},
challenge = 9,
]
\DndMonsterAction{Shapechanger}
The slaad can use its action to polymorph into a Small or Medium humanoid, or back into its true form. Its statistics, other than its size, are the same in each form. Any equipment it is wearing or carrying isn't transformed. It reverts to its true form if it dies.

\DndMonsterAction{Magic Resistance}
The slaad has advantage on saving throws against spells and other magical effects.

\DndMonsterAction{Magic Weapons}
The slaad's weapon attacks are magical.

\DndMonsterAction{Regeneration}
The slaad regains 10 hit points at the start of its turn if it has at least 1 hit point.

\DndMonsterAction{Innate Spellcasting}
The slaad's innate spellcasting ability is Charisma (spell save DC 14). The slaad can innately cast the following spells, requiring no material components:
\begin{DndMonsterSpells}
  \DndInnateSpellLevel{detect magic, \textit{detect thoughts}, \textit{invisibility} (self only), \textit{mage hand}, \textit{major image}}
  \DndInnateSpellLevel[]{plane shift (self only)}
  \DndInnateSpellLevel[2]{fear, \textit{fly}, \textit{fireball}, \textit{tongues}}
\end{DndMonsterSpells}
\DndMonsterSection{Actions}
\DndMonsterAction{Multiattack}
The slaad makes three attacks: one with its bite and two with its claws or greatsword.

\DndMonsterAction{Bite (Slaad Form Only)}
\textsl{Melee Weapon Attack:} 7 to hit, reach 5 ft., one target. \textsl{Hit:} 6 (1d6 + 3) piercing damage.

\DndMonsterAction{Claws (Slaad Form Only)}
\textsl{Melee Weapon Attack:} 7 to hit, reach 5 ft., one target. \textsl{Hit:} 8 (1d10 + 3) slashing damage.

\DndMonsterAction{Greatsword}
\textsl{Melee Weapon Attack:} 7 to hit, reach 5 ft., one target. \textsl{Hit:} 10 (2d6 + 3) slashing damage.

\end{DndMonster}



\begin{DndMonster}[float=t]{Ignatius Inkblot}
\DndMonsterType{Medium aberration, lawful evil}
\DndMonsterBasics[
armor-class = {15 (\textit{breastplate})},
hit-points = {\DndDice{13d8 + 13}},
speed = {30 ft.}
]
\DndMonsterAbilityScores[
 str = 11,
 dex = 12,
 con = 12,
 int = 19,
 wis = 17,
 cha = 17
]
\DndMonsterDetails[
saving-throws = {Int +7, Wis +6, Cha +6},
skills = {Arcana +7, Deception +6, Insight +6, Perception +6, Persuasion +6, Stealth +4},
senses = {darkvision 120 ft., passive perception 16},
languages = {Deep Speech, Undercommon, telepathy 120 ft.},
challenge = 7,
]
\DndMonsterAction{Copy Warning}
This content has been copied from another creature, but the data replacement hasn't been run. Check details.

\DndMonsterAction{Magic Resistance}
The mind flayer has advantage on saving throws against spells and other magical effects.

\DndMonsterAction{Innate Spellcasting (Psionics)}
The mind flayer's innate spellcasting ability is Intelligence (spell save DC 15). It can innately cast the following spells, requiring no components:
\begin{DndMonsterSpells}
  \DndInnateSpellLevel{detect thoughts, \textit{levitate}}
  \DndInnateSpellLevel[1]{dominate monster, \textit{plane shift} (self only)}
\end{DndMonsterSpells}
\DndMonsterSection{Actions}
\DndMonsterAction{Tentacles}
\textsl{Melee Weapon Attack:} 7 to hit, reach 5 ft., one creature. \textsl{Hit:} 15 (2d10 + 4) psychic damage. If the target is Medium or smaller, it is grappled (escape DC 15) and must succeed on a DC 15 Intelligence saving throw or be stunned until this grapple ends.

\DndMonsterAction{Extract Brain}
\textsl{Melee Weapon Attack:} 7 to hit, reach 5 ft., one incapacitated humanoid grappled by the mind flayer. \textsl{Hit:} The target takes 55 (10d10) piercing damage. If this damage reduces the target to 0 hit points, the mind flayer kills the target by extracting and devouring its brain.

\DndMonsterAction{Mind Blast (Recharge 5\textendash 6)}
The mind flayer magically emits psychic energy in a 60-foot cone. Each creature in that area must succeed on a DC 15 Intelligence saving throw or take 22 (4d8 + 4) psychic damage and be stunned for 1 minute. A creature can repeat the saving throw at the end of each of its turns, ending the effect on itself on a success.

\end{DndMonster}



\begin{DndMonster}[float=t]{Imp}
\DndMonsterType{Tiny fiend (devil), lawful evil}
\DndMonsterBasics[
armor-class = {13},
hit-points = {\DndDice{3d4 + 3}},
speed = {20 ft., fly 40 ft.}
]
\DndMonsterAbilityScores[
 str = 6,
 dex = 17,
 con = 13,
 int = 11,
 wis = 12,
 cha = 14
]
\DndMonsterDetails[
skills = {Deception +4, Insight +3, Persuasion +4, Stealth +5},
damage-resistances = {cold; bludgeoning, piercing, slashing from nonmagical attacks not made with silvered weapons},
damage-immunities = {fire, poison},
condition-immunities = {poisoned},
senses = {darkvision 120 ft., passive perception 11},
languages = {Infernal, Common},
challenge = 1,
]
\DndMonsterAction{Shapechanger}
The imp can use its action to polymorph into a beast form that resembles a rat (speed 20 ft.), a raven (20 ft., fly 60 ft.), or a spider (20 ft., climb 20 ft.), or back into its true form. Its statistics are the same in each form, except for the speed changes noted. Any equipment it is wearing or carrying isn't transformed. It reverts to its true form if it dies.

\DndMonsterAction{Devil's Sight}
Magical darkness doesn't impede the imp's darkvision.

\DndMonsterAction{Magic Resistance}
The imp has advantage on saving throws against spells and other magical effects.

\DndMonsterSection{Actions}
\DndMonsterAction{Sting (Bite in Beast Form)}
\textsl{Melee Weapon Attack:} 5 to hit, reach 5 ft., one target. \textsl{Hit:} 5 (1d4 + 3) piercing damage, and the target must make a DC 11 Constitution saving throw, taking 10 (3d6) poison damage on a failed save, or half as much damage on a successful one.

\DndMonsterAction{Invisibility}
The imp magically turns invisible until it attacks, or until its concentration ends (as if concentration on a spell). Any equipment the imp wears or carries is invisible with it.

\end{DndMonster}



\begin{DndMonster}[float=t]{Mage}
\DndMonsterType{Medium humanoid (any race), A}
\DndMonsterBasics[
armor-class = {12 (15 with \textit{mage armor})},
hit-points = {\DndDice{9d8}},
speed = {30 ft.}
]
\DndMonsterAbilityScores[
 str = 9,
 dex = 14,
 con = 11,
 int = 17,
 wis = 12,
 cha = 11
]
\DndMonsterDetails[
saving-throws = {Int +6, Wis +4},
skills = {Arcana +6, History +6},
languages = {any four languages},
challenge = 6,
]
\DndMonsterAction{Spellcasting}
The mage is a 9th-level spellcaster. Its spellcasting ability is Intelligence (spell save DC 14, 6 to hit with spell attacks). The mage has the following wizard spells prepared:
\begin{DndMonsterSpells}
  \DndMonsterSpellLevel{fire bolt, \textit{light}, \textit{mage hand}, \textit{prestidigitation}}
   \DndMonsterSpellLevel[1][4]{detect magic, \textit{mage armor}, \textit{magic missile}, \textit{shield}}
   \DndMonsterSpellLevel[2][3]{misty step, \textit{suggestion}}
   \DndMonsterSpellLevel[3][3]{counterspell, \textit{fireball}, \textit{fly}}
   \DndMonsterSpellLevel[4][3]{greater invisibility, \textit{ice storm}}
   \DndMonsterSpellLevel[5][1]{cone of cold}
\end{DndMonsterSpells}
\DndMonsterSection{Actions}
\DndMonsterAction{Dagger}
\textsl{Melee or Ranged Weapon Attack:} 5 to hit, reach 5 ft. or range 20/60 ft., one target. \textsl{Hit:} 4 (1d4 + 2) piercing damage.

\end{DndMonster}



\begin{DndMonster}[float=t]{Mind Flayer}
\DndMonsterType{Medium aberration, lawful evil}
\DndMonsterBasics[
armor-class = {15 (\textit{breastplate})},
hit-points = {\DndDice{13d8 + 13}},
speed = {30 ft.}
]
\DndMonsterAbilityScores[
 str = 11,
 dex = 12,
 con = 12,
 int = 19,
 wis = 17,
 cha = 17
]
\DndMonsterDetails[
saving-throws = {Int +7, Wis +6, Cha +6},
skills = {Arcana +7, Deception +6, Insight +6, Perception +6, Persuasion +6, Stealth +4},
senses = {darkvision 120 ft., passive perception 16},
languages = {Deep Speech, Undercommon, telepathy 120 ft.},
challenge = 7,
]
\DndMonsterAction{Magic Resistance}
The mind flayer has advantage on saving throws against spells and other magical effects.

\DndMonsterAction{Innate Spellcasting (Psionics)}
The mind flayer's innate spellcasting ability is Intelligence (spell save DC 15). It can innately cast the following spells, requiring no components:
\begin{DndMonsterSpells}
  \DndInnateSpellLevel{detect thoughts, \textit{levitate}}
  \DndInnateSpellLevel[1]{dominate monster, \textit{plane shift} (self only)}
\end{DndMonsterSpells}
\DndMonsterSection{Actions}
\DndMonsterAction{Tentacles}
\textsl{Melee Weapon Attack:} 7 to hit, reach 5 ft., one creature. \textsl{Hit:} 15 (2d10 + 4) psychic damage. If the target is Medium or smaller, it is grappled (escape DC 15) and must succeed on a DC 15 Intelligence saving throw or be stunned until this grapple ends.

\DndMonsterAction{Extract Brain}
\textsl{Melee Weapon Attack:} 7 to hit, reach 5 ft., one incapacitated humanoid grappled by the mind flayer. \textsl{Hit:} The target takes 55 (10d10) piercing damage. If this damage reduces the target to 0 hit points, the mind flayer kills the target by extracting and devouring its brain.

\DndMonsterAction{Mind Blast (Recharge 5\textendash 6)}
The mind flayer magically emits psychic energy in a 60-foot cone. Each creature in that area must succeed on a DC 15 Intelligence saving throw or take 22 (4d8 + 4) psychic damage and be stunned for 1 minute. A creature can repeat the saving throw at the end of each of its turns, ending the effect on itself on a success.

\end{DndMonster}



\begin{DndMonster}[float=t]{Monodrone}
\DndMonsterType{Medium construct, lawful neutral}
\DndMonsterBasics[
armor-class = {15 (natural armor)},
hit-points = {\DndDice{1d8 + 1}},
speed = {30 ft., fly 30 ft.}
]
\DndMonsterAbilityScores[
 str = 10,
 dex = 13,
 con = 12,
 int = 4,
 wis = 10,
 cha = 5
]
\DndMonsterDetails[
senses = {truesight 120 ft., passive perception 10},
languages = {Modron},
challenge = 1/8,
]
\DndMonsterAction{Axiomatic Mind}
The monodrone can't be compelled to act in a manner contrary to its nature or its instructions.

\DndMonsterAction{Disintegration}
If the monodrone dies, its body disintegrates into dust, leaving behind its weapons and anything else it was carrying.

\DndMonsterSection{Actions}
\DndMonsterAction{Dagger}
\textsl{Melee Weapon Attack:} 3 to hit, reach 5 ft., one target. \textsl{Hit:} 3 (1d4 + 1) piercing damage.

\DndMonsterAction{Javelin}
\textsl{Melee or Ranged Weapon Attack:} 2 to hit, reach 5 ft. or range 30/120 ft., one target. \textsl{Hit:} 3 (1d6) piercing damage.

\end{DndMonster}



\begin{DndMonster}[float=t]{Nixylanna Vidorant}
\DndMonsterType{Medium humanoid (any race), any non-good alignment}
\DndMonsterBasics[
armor-class = {15 (studded leather armor)},
hit-points = {\DndDice{12d8 + 24}},
speed = {30 ft.}
]
\DndMonsterAbilityScores[
 str = 11,
 dex = 16,
 con = 14,
 int = 13,
 wis = 11,
 cha = 10
]
\DndMonsterDetails[
saving-throws = {Dex +6, Int +4},
skills = {Acrobatics +6, Deception +3, Perception +3, Stealth +9},
damage-resistances = {poison},
languages = {Thieves' cant plus any two languages},
challenge = 8,
]
\DndMonsterAction{Copy Warning}
This content has been copied from another creature, but the data replacement hasn't been run. Check details.

\DndMonsterAction{Assassinate}
During its first turn, the assassin has advantage on attack rolls against any creature that hasn't taken a turn. Any hit the assassin scores against a surprised creature is a critical hit.

\DndMonsterAction{Evasion}
If the assassin is subjected to an effect that allows it to make a Dexterity saving throw to take only half damage, the assassin instead takes no damage if it succeeds on the saving throw, and only half damage if it fails.

\DndMonsterAction{Sneak Attack (1/Turn)}
The assassin deals an extra 14 (4d6) damage when it hits a target with a weapon attack and has advantage on the attack roll, or when the target is within 5 feet of an ally of the assassin that isn't incapacitated and the assassin doesn't have disadvantage on the attack roll.

\DndMonsterSection{Actions}
\DndMonsterAction{Multiattack}
The assassin makes two shortsword attacks.

\DndMonsterAction{Shortsword}
\textsl{Melee Weapon Attack:} 6 to hit, reach 5 ft., one target. \textsl{Hit:} 6 (1d6 + 3) piercing damage, and the target must make a DC 15 Constitution saving throw, taking 24 (7d6) poison damage on a failed save, or half as much damage on a successful one.

\DndMonsterAction{Light Crossbow}
\textsl{Ranged Weapon Attack:} 6 to hit, range 80/320 ft., one target. \textsl{Hit:} 7 (1d8 + 3) piercing damage, and the target must make a DC 15 Constitution saving throw, taking 24 (7d6) poison damage on a failed save, or half as much damage on a successful one.

\end{DndMonster}



\begin{DndMonster}[float=t]{Noble}
\DndMonsterType{Medium humanoid (any race), A}
\DndMonsterBasics[
armor-class = {15 (\textit{breastplate})},
hit-points = {\DndDice{2d8}},
speed = {30 ft.}
]
\DndMonsterAbilityScores[
 str = 11,
 dex = 12,
 con = 11,
 int = 12,
 wis = 14,
 cha = 16
]
\DndMonsterDetails[
skills = {Deception +5, Insight +4, Persuasion +5},
languages = {any two languages},
challenge = 1/8,
]
\DndMonsterSection{Actions}
\DndMonsterAction{Rapier}
\textsl{Melee Weapon Attack:} 3 to hit, reach 5 ft., one target. \textsl{Hit:} 5 (1d8 + 1) piercing damage.

\DndMonsterSection{Reactions}
\DndMonsterAction{Parry}
The noble adds 2 to its AC against one melee attack that would hit it. To do so, the noble must see the attacker and be wielding a melee weapon.

\end{DndMonster}



\begin{DndMonster}[float=t]{Nycaloth}
\DndMonsterType{Large fiend (yugoloth), neutral evil}
\DndMonsterBasics[
armor-class = {18 (natural armor)},
hit-points = {\DndDice{13d10 + 52}},
speed = {40 ft., fly 60 ft.}
]
\DndMonsterAbilityScores[
 str = 20,
 dex = 11,
 con = 19,
 int = 12,
 wis = 10,
 cha = 15
]
\DndMonsterDetails[
skills = {Intimidation +6, Perception +4, Stealth +4},
damage-resistances = {cold; fire; lightning; bludgeoning, piercing, slashing from nonmagical attacks},
damage-immunities = {acid, poison},
condition-immunities = {poisoned},
senses = {blindsight 60 ft., darkvision 60 ft., passive perception 14},
languages = {Abyssal, Infernal, telepathy 60 ft.},
challenge = 9,
]
\DndMonsterAction{Magic Resistance}
The nycaloth has advantage on saving throws against spells and other magical effects.

\DndMonsterAction{Magic Weapons}
The nycaloth's weapon attacks are magical.

\DndMonsterAction{Innate Spellcasting}
The nycaloth's innate spellcasting ability is Charisma. The nycaloth can innately cast the following spells, requiring no material components:
\begin{DndMonsterSpells}
  \DndInnateSpellLevel{darkness, \textit{detect magic}, \textit{dispel magic}, \textit{invisibility} (self only), \textit{mirror image}}
\end{DndMonsterSpells}
\DndMonsterSection{Actions}
\DndMonsterAction{Multiattack}
The nycaloth makes two melee attacks, or it makes one melee attack and teleports before or after the attack.

\DndMonsterAction{Claw}
\textsl{Melee Weapon Attack:} 9 to hit, reach 5 ft., one target. \textsl{Hit:} 12 (2d6 + 5) slashing damage. If the target is a creature, it must succeed on a DC 16 Constitution saving throw or take 5 (2d4) slashing damage at the start of each of its turns due to a fiendish wound. Each time the nycaloth hits the wounded target with this attack, the damage dealt by the wound increases by 5 (2d4). Any creature can take an action to stanch the wound with a successful DC 13 Wisdom (Medicine) check. The wound also closes if the target receives magical healing.

\DndMonsterAction{Greataxe}
\textsl{Melee Weapon Attack:} 9 to hit, reach 5 ft., one target. \textsl{Hit:} 18 (2d12 + 5) slashing damage.

\DndMonsterAction{Teleport}
The nycaloth magically teleports, along with any equipment it is wearing or carrying, up to 60 feet to an unoccupied space it can see.

\end{DndMonster}



\begin{DndMonster}[float=t]{Pentadrone}
\DndMonsterType{Large construct, lawful neutral}
\DndMonsterBasics[
armor-class = {16 (natural armor)},
hit-points = {\DndDice{5d10 + 5}},
speed = {40 ft.}
]
\DndMonsterAbilityScores[
 str = 15,
 dex = 14,
 con = 12,
 int = 10,
 wis = 10,
 cha = 13
]
\DndMonsterDetails[
skills = {Perception +4},
senses = {truesight 120 ft., passive perception 14},
languages = {Modron},
challenge = 2,
]
\DndMonsterAction{Axiomatic Mind}
The pentadrone can't be compelled to act in a manner contrary to its nature or its instructions.

\DndMonsterAction{Disintegration}
If the pentadrone dies, its body disintegrates into dust, leaving behind its weapons and anything else it was carrying.

\DndMonsterSection{Actions}
\DndMonsterAction{Multiattack}
The pentadrone makes five arm attacks.

\DndMonsterAction{Arm}
\textsl{Melee Weapon Attack:} 4 to hit, reach 5 ft., one target. \textsl{Hit:} 5 (1d6 + 2) bludgeoning damage.

\DndMonsterAction{Paralysis Gas (Recharge 5\textendash 6)}
The pentadrone exhales a 30-foot cone of gas. Each creature in that area must succeed on a DC 11 Constitution saving throw or be paralyzed for 1 minute. A creature can repeat the saving throw at the end of each of its turns, ending the effect on itself on a success.

\end{DndMonster}



\begin{DndMonster}[float=t]{Quadrone}
\DndMonsterType{Medium construct, lawful neutral}
\DndMonsterBasics[
armor-class = {16 (natural armor)},
hit-points = {\DndDice{4d8 + 4}},
speed = {30 ft., fly 30 ft.}
]
\DndMonsterAbilityScores[
 str = 12,
 dex = 14,
 con = 12,
 int = 10,
 wis = 10,
 cha = 11
]
\DndMonsterDetails[
skills = {Perception +2},
senses = {truesight 120 ft., passive perception 12},
languages = {Modron},
challenge = 1,
]
\DndMonsterAction{Axiomatic Mind}
The quadrone can't be compelled to act in a manner contrary to its nature or its instructions.

\DndMonsterAction{Disintegration}
If the quadrone dies, its body disintegrates into dust, leaving behind its weapons and anything else it was carrying.

\DndMonsterSection{Actions}
\DndMonsterAction{Multiattack}
The quadrone makes two fist attacks or four shortbow attacks.

\DndMonsterAction{Fist}
\textsl{Melee Weapon Attack:} 3 to hit, reach 5 ft., one target. \textsl{Hit:} 3 (1d4 + 1) bludgeoning damage.

\DndMonsterAction{Shortbow}
\textsl{Ranged Weapon Attack:} 4 to hit, range 80/320 ft., one target. \textsl{Hit:} 5 (1d6 + 2) piercing damage.

\end{DndMonster}



\begin{DndMonster}[float=t]{Shield Guardian}
\DndMonsterType{Large construct, U}
\DndMonsterBasics[
armor-class = {17 (natural armor)},
hit-points = {\DndDice{15d10 + 60}},
speed = {30 ft.}
]
\DndMonsterAbilityScores[
 str = 18,
 dex = 8,
 con = 18,
 int = 7,
 wis = 10,
 cha = 3
]
\DndMonsterDetails[
damage-immunities = {poison},
condition-immunities = {charmed, exhaustion, frightened, paralyzed, poisoned},
senses = {blindsight 10 ft., darkvision 60 ft., passive perception 10},
languages = {understands commands given in any language but can't speak},
challenge = 7,
]
\DndMonsterAction{Bound}
The shield guardian is magically bound to an amulet. As long as the guardian and its amulet are on the same plane of existence, the amulet's wearer can telepathically call the guardian to travel to it, and the guardian knows the distance and direction to the amulet. If the guardian is within 60 feet of the amulet's wearer, half of any damage the wearer takes (rounded up) is transferred to the guardian.

\DndMonsterAction{Regeneration}
The shield guardian regains 10 hit points at the start of its turn if it has at least 1 hit point.

\DndMonsterAction{Spell Storing}
A spellcaster who wears the shield guardian's amulet can cause the guardian to store one spell of 4th level or lower. To do so, the wearer must cast the spell on the guardian. The spell has no effect but is stored within the guardian. When commanded to do so by the wearer or when a situation arises that was predefined by the spellcaster, the guardian casts the stored spell with any parameters set by the original caster, requiring no components. When the spell is cast or a new spell is stored, any previously stored spell is lost.

\DndMonsterSection{Actions}
\DndMonsterAction{Multiattack}
The guardian makes two fist attacks.

\DndMonsterAction{Fist}
\textsl{Melee Weapon Attack:} 7 to hit, reach 5 ft., one target. \textsl{Hit:} 11 (2d6 + 4) bludgeoning damage.

\DndMonsterSection{Reactions}
\DndMonsterAction{Shield}
When a creature makes an attack against the wearer of the guardian's amulet, the guardian grants a +2 bonus to the wearer's AC if the guardian is within 5 feet of the wearer.

\end{DndMonster}



\begin{DndMonster}[float*=b,width=\textwidth + 8pt]{Solar}
\begin{multicols}{2}
\DndMonsterType{Large celestial, lawful good}
\DndMonsterBasics[
armor-class = {21 (natural armor)},
hit-points = {\DndDice{18d10 + 144}},
speed = {50 ft., fly 150 ft.}
]
\DndMonsterAbilityScores[
 str = 26,
 dex = 22,
 con = 26,
 int = 25,
 wis = 25,
 cha = 30
]
\DndMonsterDetails[
saving-throws = {Int +14, Wis +14, Cha +17},
skills = {Perception +14},
damage-resistances = {radiant; bludgeoning, piercing, slashing from nonmagical attacks},
damage-immunities = {necrotic, poison},
condition-immunities = {charmed, exhaustion, frightened, poisoned},
senses = {truesight 120 ft., passive perception 24},
languages = {all, telepathy 120 ft.},
challenge = 21,
]
\DndMonsterAction{Angelic Weapons}
The solar's weapon attacks are magical. When the solar hits with any weapon, the weapon deals an extra 6d8 radiant damage (included in the attack).

\DndMonsterAction{Divine Awareness}
The solar knows if it hears a lie.

\DndMonsterAction{Magic Resistance}
The solar has advantage on saving throws against spells and other magical effects.

\DndMonsterAction{Innate Spellcasting}
The solar's spellcasting ability is Charisma (spell save DC 25). It can innately cast the following spells, requiring no material components:
\begin{DndMonsterSpells}
  \DndInnateSpellLevel{detect evil and good, \textit{invisibility} (self only)}
  \DndInnateSpellLevel[3]{blade barrier, \textit{dispel evil and good}, \textit{resurrection}}
  \DndInnateSpellLevel[1]{commune, \textit{control weather}}
\end{DndMonsterSpells}
\DndMonsterSection{Actions}
\DndMonsterAction{Multiattack}
The solar makes two greatsword attacks.

\DndMonsterAction{Greatsword}
\textsl{Melee Weapon Attack:} 15 to hit, reach 5 ft., one target. \textsl{Hit:} 22 (4d6 + 8) slashing damage plus 27 (6d8) radiant damage.

\DndMonsterAction{Slaying Longbow}
\textsl{Ranged Weapon Attack:} 13 to hit, range 150/600 ft., one target. \textsl{Hit:} 15 (2d8 + 6) piercing damage plus 27 (6d8) radiant damage. If the target is a creature that has 100 hit points or fewer, it must succeed on a DC 15 Constitution saving throw or die.

\DndMonsterAction{Flying Sword}
The solar releases its greatsword to hover magically in an unoccupied space within 5 feet of it. If the solar can see the sword, the solar can mentally command it as a bonus action to fly up to 50 feet and either make one attack against a target or return to the solar's hands. If the hovering sword is targeted by any effect, the solar is considered to be holding it. The hovering sword falls if the solar dies.

\DndMonsterAction{Healing Touch (4/Day)}
The solar touches another creature. The target magically regains 40 (8d8 + 4) hit points and is freed from any curse, disease, poison, blindness, or deafness.

\DndMonsterSection{Legendary Actions}
Solar can take 3 legendary actions, choosing from the options below. Only one legendary action can be used at a time and only at the end of another creature's turn. Solar regains spent legendary actions at the start of its turn.

\begin{DndMonsterLegendaryActions}
\DndMonsterLegendaryAction{Teleport}{The solar magically teleports, along with any equipment it is wearing or carrying, up to 120 feet to an unoccupied space it can see.
}
\DndMonsterLegendaryAction{Searing Burst (Costs 2 Actions)}{The solar emits magical, divine energy. Each creature of its choice in a 10-foot radius must make a DC 23 Dexterity saving throw, taking 14 (4d6) fire damage plus 14 (4d6) radiant damage on a failed save, or half as much damage on a successful one.
}
\DndMonsterLegendaryAction{Blinding Gaze (Costs 3 Actions)}{The solar targets one creature it can see within 30 feet of it. If the target can see it, the target must succeed on a DC 15 Constitution saving throw or be blinded until magic such as the \textit{lesser restoration} spell removes the blindness.
}
\end{DndMonsterLegendaryActions}
\end{multicols}
\end{DndMonster}


\clearpage

\begin{DndMonster}[float=t]{The Stranger}
\DndMonsterType{Medium humanoid (any race), A}
\DndMonsterBasics[
armor-class = {12 (15 with \textit{mage armor})},
hit-points = {\DndDice{18d8 + 18}},
speed = {30 ft.}
]
\DndMonsterAbilityScores[
 str = 10,
 dex = 14,
 con = 12,
 int = 20,
 wis = 15,
 cha = 16
]
\DndMonsterDetails[
saving-throws = {Int +9, Wis +6},
skills = {Arcana +13, History +13},
damage-resistances = {damage from spells; bludgeoning, piercing, slashing (from stoneskin)},
languages = {any six languages},
challenge = 12,
]
\DndMonsterAction{Copy Warning}
This content has been copied from another creature, but the data replacement hasn't been run. Check details.

\DndMonsterAction{Magic Resistance}
The archmage has advantage on saving throws against spells and other magical effects.

\DndMonsterAction{Spellcasting}
The archmage is an 18th-level spellcaster. Its spellcasting ability is Intelligence (spell save DC 17, 9 to hit with spell attacks). The archmage can cast \textit{disguise self} and \textit{invisibility} at will and has the following wizard spells prepared:
\begin{DndMonsterSpells}
  \DndInnateSpellLevel{disguise self, \textit{invisibility}}
  \DndMonsterSpellLevel{fire bolt, \textit{light}, \textit{mage hand}, \textit{prestidigitation}, \textit{shocking grasp}}
   \DndMonsterSpellLevel[1][4]{detect magic, \textit{identify}, \textit{mage armor}*, \textit{magic missile}}
   \DndMonsterSpellLevel[2][3]{detect thoughts, \textit{mirror image}, \textit{misty step}}
   \DndMonsterSpellLevel[3][3]{counterspell, \textit{fly}, \textit{lightning bolt}}
   \DndMonsterSpellLevel[4][3]{banishment, \textit{fire shield}, \textit{stoneskin}*}
   \DndMonsterSpellLevel[5][3]{cone of cold, \textit{scrying}, \textit{wall of force}}
   \DndMonsterSpellLevel[6][1]{globe of invulnerability}
   \DndMonsterSpellLevel[7][1]{teleport}
   \DndMonsterSpellLevel[8][1]{mind blank*}
\end{DndMonsterSpells}
\DndMonsterSection{Actions}
\DndMonsterAction{Dagger}
\textsl{Melee or Ranged Weapon Attack:} 6 to hit, reach 5 ft. or range 20/60 ft., one target. \textsl{Hit:} 4 (1d4 + 2) piercing damage.

\end{DndMonster}



\begin{DndMonster}[float=t]{Tridrone}
\DndMonsterType{Medium construct, lawful neutral}
\DndMonsterBasics[
armor-class = {15 (natural armor)},
hit-points = {\DndDice{3d8 + 3}},
speed = {30 ft.}
]
\DndMonsterAbilityScores[
 str = 12,
 dex = 13,
 con = 12,
 int = 9,
 wis = 10,
 cha = 9
]
\DndMonsterDetails[
senses = {truesight 120 ft., passive perception 10},
languages = {Modron},
challenge = 1/2,
]
\DndMonsterAction{Axiomatic Mind}
The tridrone can't be compelled to act in a manner contrary to its nature or its instructions.

\DndMonsterAction{Disintegration}
If the tridrone dies, its body disintegrates into dust, leaving behind its weapons and anything else it was carrying.

\DndMonsterSection{Actions}
\DndMonsterAction{Multiattack}
The tridrone makes three fist attacks or three javelin attacks.

\DndMonsterAction{Fist}
\textsl{Melee Weapon Attack:} 3 to hit, reach 5 ft., one target. \textsl{Hit:} 3 (1d4 + 1) bludgeoning damage.

\DndMonsterAction{Javelin}
\textsl{Melee or Ranged Weapon Attack:} 3 to hit, reach 5 ft. or range 30/120 ft., one target. \textsl{Hit:} 4 (1d6 + 1) piercing damage.

\end{DndMonster}


\end{document}
